\documentclass[UTF8,a4paper,zihao=-4,fontset=macnew]{ctexart}

\usepackage[
  top=2.5cm,
  right=2.5cm,
  bottom=2.5cm,
  left=3cm,
  headheight=15pt,
  headsep=0.45cm,
  footskip=0.72cm
]{geometry}
\usepackage{fontspec}
\usepackage{graphicx}
\usepackage{xcolor}
\usepackage{array}
\usepackage{tabularx}
\usepackage{makecell}
\usepackage{multirow}
\usepackage{enumitem}
\usepackage{fancyhdr}
\usepackage{listings}
\usepackage{hyperref}
\usepackage{amsmath}
\usepackage{anyfontsize}
\usepackage{microtype}
\usepackage{float}
\usepackage{booktabs}
\usepackage{longtable}
\usepackage[font=small]{caption}
\usepackage[font=footnotesize]{subcaption}
\usepackage{tikz}
\usetikzlibrary{arrows.meta, positioning, shapes.geometric, fit, backgrounds, calc, matrix}
\tikzset{
  blk/.style ={rectangle, rounded corners=2pt, draw, align=center, inner sep=4pt, font=\footnotesize},
  io/.style  ={rectangle, draw, align=center, inner sep=4pt, font=\footnotesize, fill=black!4},
  dec/.style ={diamond, draw, aspect=2.2, align=center, inner sep=1pt, font=\footnotesize},
  lyr/.style ={rectangle, draw, rounded corners=2pt, align=center, font=\footnotesize, fill=black!3},
  grp/.style ={rectangle, rounded corners=3pt, draw=black!35, dashed, inner sep=6pt},
  fl/.style  ={-{Latex[length=2mm]}, >=Latex},
  fl2/.style ={{Latex[length=2mm]}-{Latex[length=2mm]}},
  lbl/.style ={font=\scriptsize, fill=white, inner sep=1pt},
}

\setmainfont{Times New Roman}
\setsansfont{Arial}
\setmonofont{Menlo}

\hypersetup{hidelinks}
\setlength{\parindent}{2em}
\setlength{\parskip}{0pt}
\setlength{\emergencystretch}{3em}
\tolerance=2000
\hbadness=10000
\renewcommand{\arraystretch}{1.0}
\definecolor{deepblue}{RGB}{0,0,128}

\newcolumntype{C}[1]{>{\centering\arraybackslash}p{#1}}
\newcolumntype{L}[1]{>{\raggedright\arraybackslash}p{#1}}
\newcommand{\rednote}[1]{{\color{red}#1}}
\newcommand{\bluenote}[1]{{\color{blue}#1}}
\newcommand{\deepbluenote}[1]{{\color{deepblue}#1}}
\newcommand{\notepara}[2]{%
  {\setlength{\leftskip}{2em}\setlength{\parindent}{0pt}\color{#1}#2\par}%
}
\newcommand{\bluepara}[1]{\notepara{blue}{#1}}
\newcommand{\redpara}[1]{\notepara{red}{#1}}
\newcommand{\deepbluepara}[1]{\notepara{deepblue}{#1}}
\newcommand{\blankline}[1]{\makebox[#1]{\hrulefill}}
% 非数学模式下划线：避免 \underline 进入数学模式触发 cmex 字号替换
\newcommand{\ulfield}[2]{\vtop{\hbox{\makebox[#1][c]{#2}}\kern1.5pt\hrule height 0.4pt}}
\newcommand{\coverrow}[4]{%
  \noindent\makebox[2.0cm][l]{#1}\ulfield{#2}{#3}\makebox[#4]{}\\[0.42cm]%
}

\fancypagestyle{reporthead}{
  \fancyhf{}
  \fancyhead[L]{\zihao{-5}重庆理工大学综合课程设计报告}
  \fancyhead[R]{\zihao{-5}物联网工程专业-综合课程设计I-桌面用电监控}
  \renewcommand{\headrulewidth}{0.4pt}
  \renewcommand{\footrulewidth}{0pt}
}

\fancypagestyle{reportbody}{
  \fancyhf{}
  \fancyhead[L]{\zihao{-5}重庆理工大学综合课程设计报告}
  \fancyhead[R]{\zihao{-5}物联网工程专业-综合课程设计I-桌面用电监控}
  \fancyfoot[C]{\zihao{-5}\thepage}
  \renewcommand{\headrulewidth}{0.4pt}
  \renewcommand{\footrulewidth}{0pt}
}

\lstdefinelanguage{Rust}{
  keywords={fn,let,mut,pub,extern,crate,use,mod,struct,enum,impl,trait,type,const,static,if,else,while,for,loop,match,return,break,continue,as,in,ref,self,Self,super,where,unsafe,async,await,move,no_mangle,no_std,repr},
  sensitive=true,
  morecomment=[l]{//},
  morecomment=[s]{/*}{*/},
  morestring=[b]",
  morestring=[b]',
}

\lstdefinestyle{appendixcode}{
  basicstyle=\ttfamily\zihao{5},
  commentstyle=\color{black!55},
  breaklines=true,
  columns=fullflexible,
  keepspaces=true,
  showstringspaces=false,
  tabsize=4,
  xleftmargin=2em,
  aboveskip=0.5em,
  belowskip=0.5em
}
\lstdefinestyle{serialterm}{
  basicstyle=\ttfamily\zihao{-5},
  frame=single,
  framesep=4pt,
  breaklines=true,
  breakindent=2em,
  columns=fullflexible,
  keepspaces=true,
  xleftmargin=1em,
  xrightmargin=1em,
  aboveskip=0.5em,
  belowskip=0.5em
}

\newcommand{\tocline}[2]{%
  \noindent\makebox[\textwidth][l]{\heiti\bfseries #1\dotfill #2}\par\vspace{0.05cm}%
}
\newcommand{\tocsubline}[2]{%
  \noindent\hspace{0.72cm}\makebox[\dimexpr\textwidth-0.72cm\relax][l]{\heiti\bfseries #1\dotfill #2}\par\vspace{0.05cm}%
}
\newcommand{\sectiontitle}[2]{%
  \vspace*{-0.32cm}%
  {\noindent\heiti\bfseries\fontsize{18pt}{20pt}\selectfont #1.\quad #2\par}
  \vspace{0.50cm}
}
\newcommand{\subsectiontitle}[2]{%
  \vspace{0.34cm}
  {\noindent\heiti\bfseries\fontsize{16pt}{18pt}\selectfont #1.\quad #2\par}
  \vspace{0.54cm}
}
\newcommand{\itemline}[2]{%
  {\setlength{\leftskip}{2em}\setlength{\parindent}{0pt}%
  \noindent(#1)\quad #2\par}%
}
\newcommand{\blueitemline}[2]{%
  {\setlength{\leftskip}{2em}\setlength{\parindent}{0pt}\color{blue}%
  \noindent(#1)\quad #2\par}%
}
\newcommand{\spacedblueitemline}[2]{%
  \blueitemline{#1}{#2}\vspace{0.88cm}%
}
\newcommand{\lineblueitemline}[2]{%
  \blueitemline{#1}{#2}\vspace{0.18cm}%
}

\begin{document}

% Page 1: cover.
\thispagestyle{empty}
\zihao{5}
\noindent
{\setlength{\tabcolsep}{2pt}
\begin{tabular}{|C{1.65cm}|C{3.80cm}|C{3.80cm}|}
\hline
\multicolumn{3}{|c|}{\heiti 总成绩（30 分）}\\
\hline
\multirowcell{2}{\heiti 项目} & \heiti 通讯协议实现及分析 & \heiti 报告质量及团队协作\\
 & \heiti 15 分 & \heiti 15 分\\
\hline
\heiti 得分 & & \\
\hline
\heiti 评分人 & \multicolumn{2}{c|}{}\\
\hline
\end{tabular}
}

\vspace{1.20cm}
\begin{center}
  
\includegraphics[width=10.4cm]{figures/cqut-logo-wordmark.png}

  \vspace{1.35cm}
  {\heiti\bfseries\zihao{-0} 综合课程设计I 报告\par}

  \vspace{1.05cm}
  {\heiti\zihao{3}题目：\quad \ulfield{8.0cm}{桌面用电监控设计实现}\par}

  \vspace{1.65cm}
  \begin{minipage}{9.8cm}
    \zihao{4}\raggedright
    \coverrow{二级学院}{7.0cm}{计算机科学与工程学院}{0pt}
    \coverrow{专\qquad 业}{7.0cm}{物联网工程}{0pt}
    \coverrow{班\qquad 级}{7.0cm}{}{0pt}
    \noindent 学生姓名\ulfield{3.0cm}{}学号\ulfield{2.9cm}{}\\[0.42cm]
    \coverrow{指导教师}{7.0cm}{}{0pt}
    \coverrow{时\qquad 间}{7.0cm}{2026.6}{0pt}
  \end{minipage}
\end{center}

\clearpage

% Page 3: abstract.
\pagenumbering{Roman}
\setcounter{page}{1}
\thispagestyle{empty}
\setlength{\parindent}{2em}
\phantomsection\label{sec:abstract}
\vspace*{\fill}
\begin{center}
\begin{minipage}{0.85\textwidth}
\begin{center}
{\heiti\zihao{4} 摘\quad 要}
\end{center}
\vspace{0.45cm}
\zihao{-4}

本设计基于 STM32F103RCT6（Cortex-M3）微控制器实现了一套桌面用电监控系统。系统不依赖操作系统，以裸机协作方式调度，软件按模块划分：用 PWM 输出测试频率，用定时器输入捕获测量频率，由双通道 DAC 产生正弦波，对三通道 ADC 交流信号采样并计算电参量，再经 TFT LCD 多级菜单显示、通过 USART 二进制帧协议与上位机通信。

硬件方面，系统以 MiniSTM32 开发板为核心，利用 PA8/TIM1 输出 50\% 占空比可调频率方波，PA1/TIM2 输入捕获测量外部信号频率，PA4/PA5 DAC 双通道经 DMA 循环输出 128 点正弦波表，PC0/PC3/PC2 三通道 ADC 在 TIM3 触发下以 6400 Hz 采样率完成电压、电流和漏电电流的定时顺序扫描采样，三路信号共用同一触发节拍，通道间仅有很小的转换时差；对 50 Hz 基波而言，该时差对计算结果影响很小。采集数据经 DMA 搬运后由整型算法库计算有效值、有功功率、视在功率和功率因数。通信协议采用两层设计：传输层以 A5 5A 帧头、CRC-16/Modbus 校验保证可靠性；命令层在帧内承载 ASCII 文本指令，支持 HELP、STATUS?、REPORT ON/OFF、PWM SET、DAC SET 等命令。LCD 显示采用一个主菜单与四个功能页面组成的多页面状态机，串口与 LCD 共用统一监控快照结构，保证两端显示数据来源一致。

经调试验证，系统频率测量结果以 0.01 Hz 为显示分辨率输出，PWM 软件设定范围为 1--100000 Hz，DAC 输出频率 1--1000 Hz 可调。ADC 采样、电参量计算、LCD 显示和串口通信在本次测试条件下运行稳定，能够完成课程设计要求的主要功能。

\vspace{1.55cm}
\noindent{\heiti 关键词：}STM32F103；ADC 采样；二进制帧协议；整型电参量计算
\end{minipage}
\end{center}
\vfill
\clearpage

% Page 6: manual table of contents.
\thispagestyle{reporthead}
\vspace*{0.45cm}
\begin{center}
{\heiti\zihao{3} 目\quad 录}
\end{center}
\vspace{0.58cm}
\zihao{5}
\setlength{\parindent}{0pt}
\tocline{摘\quad 要}{\pageref{sec:abstract}}
\tocline{1. 课程设计的目的及要求}{\pageref{sec:purpose}}
\tocsubline{1.1. 课程设计目的}{\pageref{sec:purpose-goal}}
\tocsubline{1.2. 课程设计要求}{\pageref{sec:purpose-req}}
\tocline{2. 系统分析}{\pageref{sec:analysis}}
\tocsubline{2.1. 功能分析}{\pageref{sec:analysis-function}}
\tocsubline{2.2. 性能需求分析}{\pageref{sec:analysis-performance}}
\tocsubline{2.3. 算法分析}{\pageref{sec:analysis-algorithm}}
\tocline{3. 系统设计}{\pageref{sec:design}}
\tocsubline{3.1. 系统整体设计}{\pageref{sec:design-system}}
\tocsubline{3.2. 硬件设计}{\pageref{sec:design-hardware}}
\tocsubline{3.3. 软件设计}{\pageref{sec:design-software}}
\tocline{4. 频率测量}{\pageref{sec:frequency}}
\tocsubline{4.1. 测试频率生成}{\pageref{sec:frequency-generate}}
\tocsubline{4.2. 频率测量}{\pageref{sec:frequency-measure}}
\tocsubline{4.3. 通讯协议}{\pageref{sec:frequency-protocol}}
\tocline{5. 交流 ADC 采样}{\pageref{sec:adc}}
\tocsubline{5.1. 测试电压电流生成与实际信号来源}{\pageref{sec:adc-generate}}
\tocsubline{5.2. 信号采集}{\pageref{sec:adc-sampling}}
\tocsubline{5.3. 通信协议}{\pageref{sec:adc-protocol}}
\tocsubline{5.4. 波形曲线绘制}{\pageref{sec:adc-waveform}}
\tocline{6. 电参量计算与分析}{\pageref{sec:calculation}}
\tocsubline{6.1. 交流电参量计算方法}{\pageref{sec:calculation-method}}
\tocsubline{6.2. 程序实现}{\pageref{sec:calculation-program}}
\tocsubline{6.3. 通信设计}{\pageref{sec:calculation-comm}}
\tocline{7. 系统调试}{\pageref{sec:debug}}
\tocsubline{7.1. 通信协议分析}{\pageref{sec:debug-protocol}}
\tocsubline{7.2. 频率调试}{\pageref{sec:debug-frequency}}
\tocsubline{7.3. AD 采样调试}{\pageref{sec:debug-ad}}
\tocsubline{7.4. 电参量调试}{\pageref{sec:debug-calculation}}
\tocsubline{7.5. 联机调试}{\pageref{sec:debug-online}}
\tocline{8. 总结与展望}{\pageref{sec:summary}}

\clearpage

\pagestyle{reportbody}
\pagenumbering{arabic}
\zihao{5}
\setlength{\parindent}{2em}
\setlength{\baselineskip}{20pt}

% ==============================
% 第 1 章：课程设计的目的及要求
% ==============================
\phantomsection\label{sec:purpose}
\vspace*{0.18cm}
\sectiontitle{1}{课程设计的目的及要求}

本课程设计以桌面用电监控为主题，基于 STM32F103RCT6 微控制器开发板，设计并实现一套能生成与采集信号、计算电参量、本地显示并远程通信的嵌入式监控系统。设计过程覆盖硬件接口配置、外设驱动、通信协议、算法实现与人机交互，目的是通过完整的系统级实践锻炼嵌入式设计能力。

\phantomsection\label{sec:purpose-goal}
\subsectiontitle{1.1}{课程设计目的}
\itemline{1}{了解并熟悉常用电子元器件工作原理和功能特性。}
\itemline{2}{信号感知电路数据：掌握交流电压、电流检测电路，继电器控制方法。}
\itemline{3}{开发平台及工具使用：熟练运用 Keil MDK 软件进行单片机的 C 语言编程。}
\itemline{4}{通信及通信协议：掌握串行通讯简单协议的程序设计思想和实现。}
\itemline{5}{频率信号生成与感知：掌握基于 PWM 的频率波形生成与基于定时器捕捉功能的频率测量的程序设计思想和实现。}
\itemline{6}{正弦波信号生成：掌握基于 DMA 传送的 DA 双路带幅值、相位调整的正弦波生成方法与实现。}
\itemline{7}{交流信号采集：掌握基于定时启动、DMA 传送的 AD 转换程序设计与实现。}
\itemline{8}{交流波形绘制：掌握交流电压、电流 LCD 屏幕绘制程序设计。}
\itemline{9}{信号分析：掌握交流电参量计算算法与程序实现。}
\itemline{10}{培养嵌入式系统设计、分析及展示能力。}

\phantomsection\label{sec:purpose-req}
\subsectiontitle{1.2}{课程设计要求}
\itemline{1}{仔细阅读课程设计指导书，分析桌面用电电路原理；}
\itemline{2}{利用 AI 帮助分析设计原理、步骤、实现；}
\itemline{3}{设计频率测试信号；}
\itemline{4}{设计修改频率测试信号参数通信协议；}
\itemline{5}{设计频率测量程序；}
\itemline{6}{设计用于交流电压电流模拟测试信号；}
\itemline{7}{设计修改交流电压电流模拟测试信号参数通信协议；}
\itemline{8}{设计交流采样程序；}
\itemline{9}{曲线绘制；}
\itemline{10}{设计电参量计算算法及程序；}
\itemline{11}{串行通讯协议设计及串口助手调试显示；}
\itemline{12}{系统联调；}
\itemline{13}{系统答辩；}
\itemline{14}{系统设计过程中相互交流、团队协作及其总结；}
\itemline{15}{设计总结。}

\clearpage

% ==============================
% 第 2 章：系统分析
% ==============================
\phantomsection\label{sec:analysis}
\sectiontitle{2}{系统分析}

本系统的核心任务是对桌面用电环境中的电信号进行生成、采集、计算和展示。系统围绕 STM32F103RCT6 微控制器，综合运用定时器、DAC、ADC、DMA 和 USART 等片上外设，实现信号生成、频率测量、交流采样、电参量计算和串口通信等功能。本章分析系统的功能需求、性能需求与关键算法。

\phantomsection\label{sec:analysis-function}
\subsectiontitle{2.1}{功能分析}

\itemline{1}{频率测试信号生成}

测试频率源由 PA8 的 TIM1\_CH1 通道产生，输出 50\% 占空比方波。默认 50 Hz，可经 LCD 菜单按键或串口命令 PWM SET 在 1--100000 Hz 内调整。定时器工作于向上计数模式，按预分频器（PSC）和自动重装载值（ARR）覆盖宽频率范围，CCR 取 ARR 的一半，得到 50\% 占空比。

\itemline{2}{频率测量}

频率测量利用 PA1 的 TIM2\_CH2 输入捕获完成。TIM2 以 1 MHz 时钟运行（72 MHz / 72），计数器每加 1 对应 1 $\mu$s。相邻上升沿的捕获时间差即为信号周期，再由 $f = 100{,}000{,}000 / T_{\mu s}$ 换算为 Hz$\times$100 格式的整数频率值。16 位计数器的溢出中断把时间高位扩展到 32 位，以支持低频信号测量。

\itemline{3}{正弦波测试信号生成}

正弦波测试信号由 PA4/PA5 的 DAC 双通道生成。系统预先算好 128 点满幅正弦查找表（0--4095，12 位分辨率），运行时按用户设定的幅度和相位缩放后存入 DMA 缓冲区。TIM6 TRGO 以固定采样率触发 DMA 循环搬运，把波形写入 DAC 数据寄存器。支持 SINGLE（单通道正弦 + 单通道中点）与 DUAL（双通道正弦，可设相位差）两种模式，频率 1--1000 Hz、幅度 0--2047 码值可调。

\itemline{4}{桌面用电信号采集与 miniIO 信号处理电路}

系统实际用电信号由 miniIO 副板完成前级检测与处理。miniIO 副板包含交流电压检测、交流电流检测、漏电电流检测、市电频率整形、继电器控制、蜂鸣器控制和 MAX7219 数码管接口等模块。本系统主要使用其中的交流电压、电流、漏电电流和频率检测部分。

交流电压检测链路为：市电电压信号经副板电压检测电路隔离和缩放后，进入运算放大器信号处理电路，并叠加 REF-1.65V 直流偏置，最终输出 VL 信号。VL 为以 1.65 V 为中心的交流采样信号，输入 STM32 的 PC0/ADC1\_IN10 通道。交流电流检测链路为：负载电流经电流互感器转换为小电流信号，经采样电阻和运放处理后输出 iL，输入 STM32 的 PC3/ADC1\_IN13 通道。漏电电流检测链路为：火线和零线共同穿过漏电互感器，当两线电流不平衡时产生差动检测信号，经运放处理后输出 iLK，输入 STM32 的 PC2/ADC1\_IN12 通道。

STM32F103 的 ADC 输入范围为 0--3.3 V，miniIO 副板用 1.65 V 基准偏置把交流信号抬升到该范围内。软件计算电参量时先扣除直流偏置，再求电压、电流和漏电电流的有效值、有功功率、视在功率和功率因数。副板输出的 fr 信号为市电交流频率整形后的方波信号，输入 PA1/TIM2\_CH2，用于定时器输入捕获测量市电频率。

ADC1 配置为扫描模式，由 TIM3 TRGO 以 6400 Hz 触发（50 Hz $\times$ 128 点），每次触发完成 3 通道规则组转换。DMA1\_Channel1 以 2 倍帧长（768 半字）循环搬运，HT 和 TC 中断分别在半帧和整帧完成时进行通道数据拆分，实现连续无丢失采样。DAC 正弦波输出仅作为测试信号源，用于在不接入 miniIO 和实际负载时验证 ADC 采样和电参量计算链路，不属于实际桌面用电检测链路。

\itemline{5}{曲线绘制}

LCD 显示模块在 DA Wave 页面绘制 DAC 双通道波形预览曲线。将 128 点 DAC 样点（0--4095）线性映射到图形区域 y 坐标，采用逐点连线方式绘制。坐标轴含纵轴（幅值方向）、横轴（零线）和刻度标记，CH1 用绿色、CH2 用黄色区分显示。在 AC Measure 页面以文本形式显示电参量数值和 ADC 原始样点范围诊断信息。

\itemline{6}{信号分析}

电参量计算采用整型算法实现，算法库通过 C 接口供主控程序调用。计算内容包括：电压有效值（VRMS）、电流有效值（IRMS）、漏电电流有效值、有功功率（P）、视在功率（S）和功率因数（PF）。所有物理量使用整数表示，避免浮点运算，适合 Cortex-M3 无 FPU 的硬件特点。

\itemline{7}{通讯协议}

系统采用两层协议设计。传输层定义二进制帧结构：A5 5A 帧头 + 版本号 + 帧类型 + 序号 + 长度 + Payload + CRC-16/Modbus 校验，最大 Payload 63 字节。帧类型包括 0x01（命令）、0x81（响应）、0x82（事件/自动上报）、0x83（协议错误）。命令层在帧 Payload 中承载 ASCII 文本指令，支持 HELP、STATUS?、REPORT?、REPORT ON/OFF、PWM SET、DAC SET MODE/FREQ/AMP/PHASE、CAL ZERO 等命令，不区分大小写。

\phantomsection\label{sec:analysis-performance}
\subsectiontitle{2.2}{性能需求分析}

\itemline{1}{硬件电路需求分析。}

STM32F103RCT6 提供 256 KB Flash 和 48 KB RAM，系统时钟 72 MHz。系统不使用堆分配，所有缓冲区均为静态数组。DAC 波形表 128 点$\times$3 组（单通道$\times$2 + 双通道交错）共约 768 字节，ADC DMA 缓冲区 768 半字（1536 字节），协议环形缓冲 256 字节，加上命令响应结构体（4$\times$256 = 1024 字节），总静态内存占用约 4 KB，远低于 48 KB 上限。ADC 时钟需 $\leq$ 14 MHz，实际取 PCLK2/6 = 12 MHz。

\itemline{2}{测量需求分析。}

频率测量分辨率由 TIM2 的 1 MHz 计数时钟决定，周期测量精度为 $\pm$1 $\mu$s。以 50 Hz 信号为例，理论周期 20000 $\mu$s，频率分辨率为 0.01 Hz（对应 $\pm$1 $\mu$s 周期误差引起的频率偏差约 0.003 Hz）。ADC 采样率 6400 Hz，每周期 128 点，满足 Nyquist 定理对 50 Hz 基波及低次谐波的采样要求。RMS 计算基于全周期采样，128 点足以收敛。

\itemline{3}{曲线及参数显示分析。}

LCD 采用 ALIENTEK TFT 屏（16 位并行接口，分辨率 240$\times$320），显示刷新周期 500 ms 或由状态变化即时触发。波形图区域约 216$\times$120 像素，128 点波形经线性插值连线后视觉平滑。电参量以文本格式显示，显示分辨率分别为电压 0.01 V、电流 0.001 A、功率 0.1 W。

\itemline{4}{调试设计分析。}

USART1 配置为 115200 bps、8N1 格式，通过 CH340 USB 转串口芯片连接上位机。协议帧超时保护 300 ms，半帧间隔超时自动重置解析器。环形缓冲区溢出检测保证高负载下不会进入错误状态。自动上报功能可按固定周期（约 1 秒）向串口推送 STATUS 快照，便于实时监控。

\phantomsection\label{sec:analysis-algorithm}
\subsectiontitle{2.3}{算法分析}

本系统各功能模块涉及的核心算法如下：

（1）PWM 频率生成算法。给定目标频率 $f$，先计算总计数值 $N = f_{clk} / f$（$f_{clk} = 72$ MHz），再由 $N$ 拆分出预分频器 $\text{PSC} = \lceil N / 65536 \rceil - 1$ 和自动重装载值 $\text{ARR} = f_{clk} / (f \cdot (\text{PSC}+1)) - 1$。PWM 比较值 CCR = (ARR+1)/2 保证 50\% 占空比。

（2）输入捕获测频算法。TIM2 以 1 MHz 运行，相邻上升沿捕获值相减得到周期 $T$（$\mu$s），频率 $f = 100{,}000{,}000 / T$（单位 Hz$\times$100）。32 位时间戳通过 16 位溢出计数器扩展，支持最长约 71 分钟的无环绕时间表示。

（3）DAC 正弦波生成算法。预计算 128 点满幅正弦查找表 $S[n] = 2048 + 2047 \cdot \sin(2\pi n / 128)$，运行时按幅度参数 $A$ 缩放：$S'[n] = 2048 + (S[n] - 2048) \times A / 2047$。双通道模式下 CH2 引入相位偏移索引 $\Delta = \phi \times 128 / 360$。TIM6 采样率 $f_s = f_{sig} \times 128$，通过调整 PSC 和 ARR 实现。

（4）ADC 采样与数据拆分算法。DMA 以 circular 模式搬运 2 倍帧长（768 半字），每帧包含 3 通道 $\times$ 128 点交错数据。HT 中断处理前半帧、TC 中断处理后半帧，各自将交错序列 V, I, ILK, V, I, ILK, \ldots 拆分为三个独立连续数组。

（5）RMS 有效值算法。对 $N$ 个采样点先减 DC 偏移 $Z$，再计算 $\text{RMS} = \sqrt{\frac{1}{N}\sum_{i=0}^{N-1} (x_i - Z)^2}$。DC 偏移由 CAL ZERO 命令手动校准或自动取算术平均值。中间平方和用 64 位整数累加避免溢出，最终由整数开平方运算求解。

（6）有功功率算法。逐点计算电压、电流去直流后的瞬时功率积 $p_n = (v_n - Z_v) \times (i_n - Z_i)$，取均值后乘组合增益得物理功率。公式为 $P \times 10 = \overline{p} \times G_v \times G_i / 10^{10}$，保留符号（正=消耗，负=回馈）。视在功率 $S = V_{rms} \times I_{rms}$，功率因数 $\text{PF} = |P| / S$，$S = 0$ 时 PF 返回 0。

\clearpage

% ==============================
% 第 3 章：系统设计
% ==============================
\phantomsection\label{sec:design}
\sectiontitle{3}{系统设计}

本章介绍系统整体架构、硬件接口分配与软件模块设计。

\phantomsection\label{sec:design-system}
\subsectiontitle{3.1}{系统整体设计}

系统以 STM32F103RCT6 微控制器为核心，无 RTOS 依赖，采用裸机协作调度架构。上电后由启动代码依次完成板级安全配置、LCD 初始化、按键输入、串口协议、PWM 输出、输入捕获、DAC 波形输出、ADC 采样、命令解析和 UI 状态机的初始化，随后进入主循环。

主循环每轮迭代按固定顺序执行以下任务：获取统一时间戳 $\rightarrow$ 串口协议解析与命令分发 $\rightarrow$ 捕获超时与上报节流 $\rightarrow$ ADC 新帧处理与电参量计算 $\rightarrow$ 自动上报判断 $\rightarrow$ 按键扫描与 UI 状态机 $\rightarrow$ 监控快照读取 $\rightarrow$ LCD 显示刷新。同一轮循环共用同一个 now\_us 时间基准，避免不同模块使用不一致的时间戳。

% [占位图]
\begin{figure}[H]
\centering
\begin{subfigure}[b]{\textwidth}
\centering
\resizebox{\linewidth}{!}{%
\begin{tikzpicture}[node distance=6mm and 9mm]
  \node[io] (mains) {市电输入\\/ 负载};
  \node[blk, right=of mains] (mini) {miniIO 副板\\信号处理\\电压/电流/漏电/频率};
  \node[blk, right=of mini] (jp1) {JP1\\排针};
  \node[blk, minimum height=16mm, right=of jp1] (mcu) {STM32F103RCT6\\ADC1\,(PC0/PC3/PC2)\\TIM2 捕获\,(PA1)\\电参量算法};
  \node[io, above right=2mm and 10mm of mcu] (lcd) {LCD 显示};
  \node[io, below right=2mm and 10mm of mcu] (uart) {USART1\\上位机};
  \node[io, below=8mm of mcu] (keys) {按键\\KEYUP / KEY1 / KEY0};
  \draw[fl] (mains)--(mini);
  \draw[fl] (mini)--(jp1);
  \draw[fl] (jp1)--(mcu);
  \draw[fl] (mcu.east)--(lcd.west);
  \draw[fl] (mcu.east)--(uart.west);
  \draw[fl] (keys)--(mcu);
\end{tikzpicture}}
\caption{实际监控链路}
\end{subfigure}

\vspace{4mm}

\begin{subfigure}[b]{\textwidth}
\centering
\resizebox{\linewidth}{!}{%
\begin{tikzpicture}[node distance=6mm and 9mm]
  \node[blk] (dac) {STM32 DAC\\PA4 / PA5};
  \node[io, right=of dac] (loop) {跳线环回};
  \node[blk, right=of loop] (adc) {STM32 ADC\\PC0 / PC3};
  \node[blk, right=of adc] (calc) {电参量算法\\(验证采样链路)};
  \node[blk, below=8mm of dac] (tim1) {TIM1 PWM\\PA8 方波输出};
  \node[io, right=18mm of tim1] (pa1) {PA1 测频\\(跳线自测)};
  \draw[fl] (dac)--(loop);
  \draw[fl] (loop)--(adc);
  \draw[fl] (adc)--(calc);
  \draw[fl] (tim1)--(pa1);
\end{tikzpicture}}
\caption{模拟测试链路（DAC 环回 + PWM 自测）}
\end{subfigure}
\caption{系统整体架构框图}
\label{fig:system-arch}
\end{figure}

\phantomsection\label{sec:design-hardware}
\subsectiontitle{3.2}{硬件设计}

系统硬件以 MiniSTM32 开发板为基础，主要外设引脚分配如表 \ref{tab:pin-alloc} 所示。

\begin{table}[H]
\centering
\zihao{-5}
\caption{主要引脚分配表}
\label{tab:pin-alloc}
\begin{tabular}{llll}
\toprule
引脚 & 外设/模式 & 功能 & 方向 \\
\midrule
PA0  & GPIO 下拉输入 & KEYUP 确认键 & 输入 \\
PA1  & TIM2\_CH2 输入捕获 & 频率测量输入 & 浮空输入 \\
PA4  & DAC\_OUT1 & 正弦波输出 CH1 & 模拟输出 \\
PA5  & DAC\_OUT2 & 正弦波输出 CH2 & 模拟输出 \\
PA8  & TIM1\_CH1 PWM & 方波测试频率输出 & 复用推挽 \\
PA9  & USART1\_TX & 串口发送 & 复用推挽 \\
PA10 & USART1\_RX & 串口接收 & 浮空输入 \\
PA15 & GPIO 上拉输入 & KEY1 上键 & 输入 \\
PC0  & ADC1\_IN10 & 电压采样 VL & 模拟输入 \\
PC2  & ADC1\_IN12 & 漏电电流采样 iLK & 模拟输入 \\
PC3  & ADC1\_IN13 & 电流采样 iL & 模拟输入 \\
PC5  & GPIO 上拉输入 & KEY0 下键 & 输入 \\
PC13 & GPIO 推挽输出 & 触摸屏片选禁用 & 输出 \\
PB0--PB15 & GPIO 推挽输出 & LCD 16 位数据总线 & 输出 \\
PC6--PC10 & GPIO 推挽输出 & LCD 控制线 (RD/WR/RS/CS/BL) & 输出 \\
PD2  & GPIO 推挽输出 & LED\_RUN 心跳指示灯 & 输出 \\
\bottomrule
\end{tabular}
\end{table}

开发板上 PC0--PC3 和 PC13 与板载触摸屏（T\_SCK/T\_PEN/T\_MISO/T\_MOSI/T\_CS）存在引脚复用。软件层面在启动时将 PC13 拉高以禁用触摸屏片选，但硬件上拉电阻仍然存在，因此弃用 PC1（ADC1\_IN11，触摸屏 T\_PEN），将电流采样改迁至 PC3（ADC1\_IN13）。PA2 和 PA3 分别拉高以禁用 W25Q64 Flash 和 SD 卡的片选，避免干扰。

JTAG 接口在启动时被禁用（仅保留 SWD），因为 PB3/PB4 作为 LCD 数据线使用。

\textbf{miniIO 副板接口与实际用电检测链路}

系统实际用电检测依赖 miniIO 副板完成前级信号处理。miniIO 副板通过 JP1 排针与 MiniSTM32 开发板连接，各信号分配如表 \ref{tab:miniio-iface} 所示。

\begin{table}[H]
\centering
\zihao{-5}
\caption{miniIO 副板与 MiniSTM32 接口分配表}
\label{tab:miniio-iface}
\begin{tabular}{llll}
\toprule
miniIO 信号 & JP1 引脚 & 信号功能 & STM32 连接 \\
\midrule
5V  & 23、24 & 副板供电 & MiniSTM32 5V \\
GND & 1、2、21、22、25、26 & 电源地 & MiniSTM32 GND \\
fr  & 13、14 & 市电频率方波 & PA1 / TIM2\_CH2 \\
iLK & 15、16 & 漏电电流检测 & PC2 / ADC1\_IN12 \\
iL  & 17、18 & 负载电流检测 & PC3 / ADC1\_IN13 \\
VL  & 19、20 & 交流电压检测 & PC0 / ADC1\_IN10 \\
beep & 9、10 & 蜂鸣器控制 & 预留（未启用） \\
relay & 11、12 & 继电器控制 & 预留（未启用） \\
CLK & 3、4 & MAX7219 时钟 & 预留（未启用） \\
CS  & 5、6 & MAX7219 片选 & 预留（未启用） \\
DIN & 7、8 & MAX7219 数据输入 & 预留（未启用） \\
\bottomrule
\end{tabular}
\end{table}

miniIO 的 VL、iL、iLK 三路模拟信号均经过前端互感器和运放处理，并叠加 1.65 V 直流偏置后送入 ADC。STM32 侧不能直接接入市电信号，而是只采集副板输出的低压模拟量。fr 信号为副板输出的低压方波，可直接送入 TIM2\_CH2 进行输入捕获。继电器和蜂鸣器接口为数字输出控制信号，本次系统中未纳入软件控制，在硬件表中作为预留接口说明。MAX7219 数码管接口（CLK/CS/DIN）同样作为后续扩展预留。

% [占位图]
\begin{figure}[H]
\centering
\resizebox{0.95\textwidth}{!}{%
\begin{tikzpicture}[node distance=4mm and 26mm]
  % miniIO 副板检测模块
  \node[blk] (vl) {P1 电压检测};
  \node[blk, below=of vl] (il) {P3 电流检测};
  \node[blk, below=of il] (ilk) {P4 漏电检测};
  \node[blk, below=of ilk] (fr) {P2 频率整形};
  \node[io, below=of fr] (pwr) {供电 / 共地};
  % MiniSTM32 引脚
  \node[blk, right=of vl]  (pc0) {PC0 / ADC1\_IN10};
  \node[blk, right=of il]  (pc3) {PC3 / ADC1\_IN13};
  \node[blk, right=of ilk] (pc2) {PC2 / ADC1\_IN12};
  \node[blk, right=of fr]  (pa1) {PA1 / TIM2\_CH2};
  \node[io, right=of pwr]  (mpwr) {MiniSTM32\\5V / GND};
  % 信号连线（JP1 排针引脚号标注）
  \draw[fl] (vl)  -- node[lbl,above]{VL · JP1\,19/20} (pc0);
  \draw[fl] (il)  -- node[lbl,above]{iL · JP1\,17/18} (pc3);
  \draw[fl] (ilk) -- node[lbl,above]{iLK · JP1\,15/16} (pc2);
  \draw[fl] (fr)  -- node[lbl,above]{fr · JP1\,13/14} (pa1);
  \draw[fl2] (mpwr) -- node[lbl,above]{5V\,(JP1\,23/24) · GND} (pwr);
  % 分组
  \begin{scope}[on background layer]
    \node[grp, fit=(vl)(il)(ilk)(fr)(pwr), label=above:{miniIO 副板}] {};
    \node[grp, fit=(pc0)(pc3)(pc2)(pa1)(mpwr), label=above:{MiniSTM32 开发板}] {};
  \end{scope}
\end{tikzpicture}}
\caption{miniIO 与 MiniSTM32 硬件连接框图}
\label{fig:miniio-connect}
\end{figure}

\phantomsection\label{sec:design-software}
\subsectiontitle{3.3}{软件设计}

软件采用分层模块化架构，自底向上分为启动层、板级支持层、驱动层、应用层和用户界面层。各层职责清晰，通过头文件接口降低层间依赖。

{\setlength{\leftskip}{2em}\setlength{\parindent}{0pt}
\textbf{分层结构：}

startup/ $\rightarrow$ Reset\_Handler $\rightarrow$ SystemInit $\rightarrow$ main $\rightarrow$ app\_init $\rightarrow$ app\_run

bsp/ $\rightarrow$ board\_init：GPIO 安全电平（PC13/PA2/PA3 拉高禁用外设）、JTAG 释放、系统滴答

drivers/ $\rightarrow$ CMSIS + STM32F10x 标准外设库（仅链接实际使用的外设）

middleware/ $\rightarrow$ ALIENTEK TFT LCD 驱动（16 位并行接口）

app/ $\rightarrow$ 应用逻辑层，包含以下独立模块：

{\setlength{\leftskip}{4em}
app\_protocol：二进制帧传输层（解析/发送/CRC/环形缓冲）

app\_command：ASCII 命令处理层（命令匹配/参数解析/响应格式化）

app\_pwm：PWM 方波输出（TIM1\_CH1）

app\_capture：输入捕获测频（TIM2\_CH2）

app\_dac：DAC 正弦波输出（TIM6 + DMA2 + DAC 双通道）

app\_adc：ADC 交流采样（TIM3 + ADC1 + DMA1 三通道扫描）

app\_display：LCD 显示渲染

app\_ui：菜单状态机（5 个页面状态：主菜单、PWM、频率测量、DA 波形、AC 参数）

app\_keys：按键消抖扫描

app\_monitor\_state：统一监控快照

rust\_algos/ $\rightarrow$ no\_std 静态库，cbindgen 生成 C 头文件，提供整型电参量计算
}
}

% [占位图]
\begin{figure}[H]
\centering
\resizebox{\textwidth}{!}{%
\begin{tikzpicture}[node distance=6mm]
  \node[lyr, minimum width=12cm] (startup) {\textbf{startup/}\quad Reset\_Handler $\to$ SystemInit $\to$ main $\to$ app\_init $\to$ app\_run};
  \node[lyr, minimum width=12cm, above=of startup] (bsp) {\textbf{bsp/}\quad board\_init：GPIO 安全电平（PC13/PA2/PA3 拉高）· JTAG 释放 · 系统滴答};
  \node[lyr, minimum width=12cm, above=of bsp] (drv) {\textbf{drivers/} CMSIS + STM32F10x 标准外设库\qquad\textbf{middleware/} ALIENTEK TFT LCD 驱动};
  \matrix (app) [above=10mm of drv, matrix of nodes, ampersand replacement=\&,
      nodes={blk, anchor=center, minimum width=22mm, minimum height=6mm, font=\scriptsize},
      row sep=2mm, column sep=2mm] {
    protocol \& command \& pwm \& capture \& dac \\
    adc \& display \& ui \& keys \& monitor\_state \\
  };
  \node[blk, right=16mm of app, align=center] (rust) {\textbf{rust\_algos/}\\no\_std 整型\\电参量算法};
  \draw[fl] (startup)--(bsp);
  \draw[fl] (bsp)--(drv);
  \draw[fl] (drv.north)--(app.south);
  \draw[fl2] (app.east) -- node[lbl,above]{C 接口} node[lbl,below]{pm\_calc\_electrical} (rust);
  \begin{scope}[on background layer]
    \node[grp, fit=(app), label={[font=\footnotesize]left:\textbf{app/}}] {};
  \end{scope}
\end{tikzpicture}}
\caption{软件模块层次结构图}
\label{fig:software-arch}
\end{figure}

\textbf{主循环流程：}主循环是系统调度的核心，每轮迭代按以下顺序执行：

（1）调用 app\_capture\_get\_time\_us() 获取统一的 32 位微秒时间戳 now\_us，本轮所有模块使用同一时间基准。

（2）调用 app\_protocol\_task(now\_us) 从环形缓冲区取出字节送入帧解析状态机，解析完成的命令帧交由 app\_command\_handle\_line 处理，返回 display\_changed 标志。

（3）调用 app\_capture\_task() 处理输入捕获超时检测和自动上报节流逻辑。

（4）调用 app\_adc\_task() 检查是否有新的 ADC 帧就绪，若有则复制到工作帧并调用算法库计算电参量。

（5）判断自动上报条件：若 REPORT ON 已启用且 take\_report 信号触发，读取监控快照并格式化 STATUS 行通过事件帧发送。

（6）调用 app\_keys\_poll(now\_us) 扫描按键消抖，再由 app\_ui\_handle\_key 驱动菜单状态机。

（7）调用 app\_monitor\_state\_read 组装最新快照，判断 display\_changed 后调用 app\_display\_task 刷新 LCD。

% [占位图]
\begin{figure}[H]
\centering
\begin{tikzpicture}[node distance=7mm and 16mm, >=Latex]
  \node[blk] (s1) {取统一时间戳 now\_us};
  \node[blk, below=of s1] (s2) {协议解析与命令分发\\app\_protocol\_task};
  \node[blk, below=of s2] (s3) {捕获超时检测 / 上报节流\\app\_capture\_task};
  \node[dec, below=of s3] (d4) {ADC\\新帧就绪?};
  \node[blk, right=of d4] (a4) {复制工作帧\\$\to$ 计算电参量};
  \node[dec, below=14mm of d4] (d5) {REPORT ON\\且 take\_report?};
  \node[blk, right=of d5] (a5) {读快照 $\to$ 格式化 STATUS\\$\to$ 事件帧上报};
  \node[blk, below=14mm of d5] (s6) {按键扫描 + UI 状态机\\app\_keys\_poll / app\_ui\_handle\_key};
  \node[dec, below=of s6] (d7) {display\\\_changed?};
  \node[blk, right=of d7] (a7) {刷新 LCD\\app\_display\_task};
  \draw[fl] (s1)--(s2);
  \draw[fl] (s2)--(s3);
  \draw[fl] (s3)--(d4);
  \draw[fl] (d4)-- node[lbl,left]{继续} (d5);
  \draw[fl] (d5)-- node[lbl,left]{继续} (s6);
  \draw[fl] (s6)--(d7);
  \draw[fl] (d4)-- node[lbl,above]{是} (a4);
  \draw[fl] (d5)-- node[lbl,above]{是} (a5);
  \draw[fl] (d7)-- node[lbl,above]{是} (a7);
  \draw[fl] (d7.west) -- node[lbl,below]{下一轮循环} ++(-25mm,0) |- (s1.west);
\end{tikzpicture}
\caption{主循环流程图}
\label{fig:main-loop}
\end{figure}

\textbf{关键设计决策：}

（1）统一监控快照（app\_monitor\_state\_t）。将频率测量、PWM 输出、DAC 配置和 ADC 电参量集中到一个结构体，LCD 显示和串口 STATUS? 命令都从同一份快照读取数据，保证两端显示一致。

（2）ISR 与主循环数据隔离。串口接收采用 256 字节环形缓冲区，ISR 只负责 push、主循环负责 pop；ADC 采用三层快照架构（ISR 双缓冲 $\rightarrow$ 主循环工作帧 $\rightarrow$ 显示帧），临界区极短（仅关中断取帧编号），避免阻塞中断。

（3）Rust/C 混合编程。电参量计算使用 Rust 实现为 no\_std 静态库，通过 cbindgen 自动生成 C 头文件，CMake 构建时自动调用 cargo build 编译。接口边界使用 \texttt{\#[repr(C)]} 和 \texttt{extern "C"} 保证 ABI 兼容。

\clearpage

% ==============================
% 第 4 章：频率测量
% ==============================
\phantomsection\label{sec:frequency}
\sectiontitle{4}{频率测量}

频率测量模块包括测试频率信号生成和外部信号频率测量两部分。测试信号由 PA8/TIM1\_CH1 输出 PWM 方波，可通过 LCD 菜单或串口命令调整频率。频率测量由 PA1/TIM2\_CH2 输入捕获实现，通过相邻上升沿时间差计算频率。两部分通过统一的二进制帧协议与上位机通信。

\phantomsection\label{sec:frequency-generate}
\subsectiontitle{4.1}{测试频率生成}

PWM 测试信号由 TIM1 通道 1（PA8）输出，配置为 PWM1 模式、向上计数、50\% 占空比。频率调整的核心是根据目标频率计算定时器的预分频器 PSC 和自动重装载值 ARR。

\textbf{参数计算方法：}定时器时钟 $f_{clk} = 72$ MHz（TIM1 挂载在 APB2 总线上，APB2 预分频为 1，定时器时钟等于系统时钟）。目标总计数值 $N = f_{clk} / f_{target}$。PSC 和 ARR 均为 16 位寄存器（最大值 65535），当 $N > 65535$ 时需要预分频：

$$\text{PSC} = \lceil N / 65536 \rceil - 1$$
$$\text{ARR} = f_{clk} / (f_{target} \times (\text{PSC} + 1)) - 1$$
$$\text{CCR} = (\text{ARR} + 1) / 2$$

以默认 50 Hz 为例：$N = 72{,}000{,}000 / 50 = 1{,}440{,}000$，$\text{PSC} = \lceil 1440000 / 65536 \rceil - 1 = 21$，$\text{ARR} = 72000000 / (50 \times 22) - 1 = 65454$，$\text{CCR} = 32727$。

频率范围限制为 1--100000 Hz，超出范围自动钳位。修改频率时通过 TIM\_PrescalerConfig（立即生效模式）和 TIM\_SetAutoreload 更新参数，并触发 Update 事件使新值立即加载。

% [占位图]
\begin{figure}[H]
\centering
\begin{tikzpicture}[node distance=6mm and 14mm, >=Latex]
  \node[blk] (n1) {设置目标频率 $f_{target}$};
  \node[blk, below=of n1] (n2) {频率钳位 1--100000\,Hz};
  \node[blk, below=of n2] (n3) {计算 $N = f_{clk}/f_{target}$};
  \node[dec, below=of n3] (d1) {$N > 65535$?};
  \node[blk, right=of d1] (a1) {$\text{PSC}=\lceil N/65536\rceil-1$};
  \node[blk, below=16mm of d1] (n4) {$\text{ARR}=f_{clk}/(f_{target}(\text{PSC}{+}1))-1$\\$\text{CCR}=(\text{ARR}{+}1)/2$};
  \node[blk, below=of n4] (n5) {配置 TIM1 寄存器\\(PSC / ARR / CCR)};
  \node[blk, below=of n5] (n6) {触发 Update 事件\\新值立即加载};
  \node[blk, below=of n6] (n7) {输出 PWM 方波\\PA8, 50\% 占空比};
  \draw[fl] (n1)--(n2);
  \draw[fl] (n2)--(n3);
  \draw[fl] (n3)--(d1);
  \draw[fl] (d1)-- node[lbl,above]{是} (a1);
  \draw[fl] (a1) |- (n4.east);
  \draw[fl] (d1)-- node[lbl,left]{否：PSC$=0$} (n4);
  \draw[fl] (n4)--(n5);
  \draw[fl] (n5)--(n6);
  \draw[fl] (n6)--(n7);
\end{tikzpicture}
\caption{PWM 频率生成流程图}
\label{fig:pwm-flow}
\end{figure}

\phantomsection\label{sec:frequency-measure}
\subsectiontitle{4.2}{频率测量}

频率测量基于 TIM2 通道 2（PA1）的输入捕获功能。TIM2 配置为 1 MHz 计数时钟（72 MHz / 72），ARR=0xFFFF，以自由运行方式计数，溢出后从 0 重新计数。输入捕获配置为上升沿触发，数字滤波器系数 0x03（4 个采样时钟一致后确认边沿）。

\textbf{测量原理：}每次上升沿触发 CC2 中断，读取 16 位捕获值 CCR2。通过 32 位时间戳扩展函数 make\_capture\_time() 将 16 位值扩展为 32 位微秒时间：高 16 位来自溢出计数器 s\_overflows，低 16 位来自 CCR2。当硬件 Update 标志已置位且捕获值较小（$<$ 32768）时，说明捕获发生在溢出之后，高位需补加 1。

相邻两次捕获的 32 位时间戳相减得到信号周期 $T$（$\mu$s），频率计算为：

$$f \times 100 = 100{,}000{,}000 / T$$

结果以 Hz$\times$100 格式存储，分辨率 0.01 Hz。

\textbf{超时处理：}主循环中 app\_capture\_task() 检测距离上次捕获超过 2 秒无边沿时，将频率清零并清除捕获标志，防止 LCD 和串口显示过期的频率值。

\textbf{上报节流：}自动上报采用双重策略——周期上报每 1 秒触发一次；频率变化时额外上报，但两次变化上报间隔不小于 200 ms 以防抖。

% [占位图]
\begin{figure}[H]
\centering
\begin{tikzpicture}[node distance=6mm and 14mm, >=Latex]
  \node[blk] (c1) {上升沿触发 CC2 中断};
  \node[blk, below=of c1] (c2) {读取 16 位捕获值 CCR2};
  \node[dec, below=of c2] (d) {Update 置位\\且 CCR2 $<$ 32768?};
  \node[blk, right=of d] (a) {高 16 位 $+1$\\(捕获在溢出之后)};
  \node[blk, below=16mm of d] (c3) {make\_capture\_time()\\高 16 位(溢出计数) : 低 16 位(CCR2)\\$\to$ 32 位微秒时间戳};
  \node[blk, below=of c3] (c4) {相邻时间戳相减\\$\to$ 周期 $T$ ($\mu$s)};
  \node[blk, below=of c4] (c5) {$f\times100 = 10^{8} / T$};
  \node[blk, below=of c5] (c6) {以 Hz$\times$100 存储\\(0.01\,Hz 分辨率)};
  \draw[fl] (c1)--(c2);
  \draw[fl] (c2)--(d);
  \draw[fl] (d)-- node[lbl,above]{是} (a);
  \draw[fl] (a) |- (c3.east);
  \draw[fl] (d)-- node[lbl,left]{否} (c3);
  \draw[fl] (c3)--(c4);
  \draw[fl] (c4)--(c5);
  \draw[fl] (c5)--(c6);
\end{tikzpicture}
\caption{输入捕获测频流程图}
\label{fig:capture-flow}
\end{figure}

\phantomsection\label{sec:frequency-protocol}
\subsectiontitle{4.3}{通讯协议}

频率测量相关的串口通信通过二进制帧协议实现，协议在本模块中的具体应用如下。

\textbf{传输层帧结构：}每帧由 SOF（A5 5A）+ VERSION（0x01）+ TYPE + SEQ + LEN（小端 16 位）+ Payload + CRC-16/Modbus（小端 16 位）组成。CRC 覆盖范围从 VERSION 到 Payload 末尾，不含 SOF 和 CRC 本身。

\textbf{命令层：}帧 Payload 承载 ASCII 文本命令。与频率相关的命令包括：

\begin{table}[H]
\centering
\zihao{-5}
\caption{频率相关串口命令}
\label{tab:freq-commands}
\begin{tabular}{lll}
\toprule
命令 & 方向 & 说明 \\
\midrule
PWM SET \textit{hz} & 主机$\rightarrow$设备 & 设置 PWM 输出频率（1--100000 Hz） \\
STATUS? & 主机$\rightarrow$设备 & 查询当前状态（含 MEAS 和 PWM 字段） \\
REPORT ON / OFF & 主机$\rightarrow$设备 & 开启/关闭自动上报 \\
OK PWM FREQ=\textit{hz}Hz & 设备$\rightarrow$主机 & PWM SET 命令响应 \\
OK STATUS MEAS=\textit{f}Hz PWM=\textit{f}Hz \ldots & 设备$\rightarrow$主机 & STATUS? 响应或自动上报 \\
\bottomrule
\end{tabular}
\end{table}

\textbf{帧解析状态机：}解析器共有 10 个状态（SOF0 $\rightarrow$ SOF1 $\rightarrow$ VER $\rightarrow$ TYPE $\rightarrow$ SEQ $\rightarrow$ LEN\_L $\rightarrow$ LEN\_H $\rightarrow$ PAYLOAD $\rightarrow$ CRC\_L $\rightarrow$ CRC\_H），逐字节驱动。SOF1 状态支持连续 A5 噪声抑制——收到非 5A 字节时，若为 A5 则保持 SOF1 等待后续 5A，否则回到 SOF0。帧超时保护 300 ms，半帧超时自动重置。

\clearpage

% ==============================
% 第 5 章：交流 ADC 采样
% ==============================
\phantomsection\label{sec:adc}
\sectiontitle{5}{交流ADC采样}

交流采样模块负责对桌面交流电信号进行定时 ADC 转换并通过 DMA 传送至内存。系统同时具备 DAC 正弦波信号生成能力，可作为测试信号源供采样验证。本章包括测试信号生成、信号采集、通信协议与波形绘制。

\phantomsection\label{sec:adc-generate}
\subsectiontitle{5.1}{测试电压电流生成与实际信号来源}

本系统存在两类交流信号来源。第一类为 DAC 生成的可控正弦波测试信号，用于在不接入 miniIO 和实际负载的情况下验证 ADC 采样、DMA 搬运、电参量计算和 LCD/串口显示链路。第二类为 miniIO 副板输出的实际桌面用电检测信号，包括 VL、iL 和 iLK 三路模拟量，以及 fr 频率方波。DAC 测试信号主要用于算法与采样链路调试，miniIO 输出信号用于系统联机后的实际用电监控。

测试电压和电流信号由 PA4/PA5 的 DAC 双通道生成，采用查找表 + DMA 循环 + 定时器触发的方案。

\textbf{正弦查找表：}预计算 128 点满幅正弦表 s\_sine\_fullscale[128]，值域 0--4095（12 位 DAC 全幅），以 2048 为零点上下对称。该表在编译期确定，运行时不做 sin() 浮点运算。

\textbf{幅度缩放：}对每个满幅样点做线性缩放，公式为：

$$S'[n] = 2048 + (S[n] - 2048) \times A / 2047$$

其中 $A$ 为用户设定的幅度码值（0--2047），结果钳位到 0--4095。

\textbf{相位偏移：}DUAL 模式下 CH2 相对 CH1 的相位偏移通过索引偏移实现：$\Delta = \phi \times 128 / 360$，CH2 读取查找表时索引为 $(n + \Delta) \bmod 128$。

\textbf{DMA 传输：}SINGLE 模式使用 DMA2\_Channel3 向 DHR12R1 搬运 16 位数据（CH1 正弦波），DMA2\_Channel4 向 DHR12R2 搬运 16 位数据（CH2 固定中点 2048）。DUAL 模式仅用 DMA2\_Channel3 向 DHR12RD 搬运 32 位数据，一次搬运同时更新两通道。DMA 配置为 circular 模式，自动循环读取 128 点缓冲区。

\textbf{触发时序：}TIM6 以 TRGO（Update 事件）触发 DAC 转换。采样率 $f_s = f_{sig} \times 128$，由 TIM6 的 PSC 和 ARR 决定。以 50 Hz 信号为例，$f_s = 6400$ Hz，$N = 72{,}000{,}000 / 6400 = 11250$，PSC = 0，ARR = 11249。

% [占位图]
\begin{figure}[H]
\centering
\resizebox{\textwidth}{!}{%
\begin{tikzpicture}[node distance=12mm, >=Latex]
  \node[blk] (t1) {128 点正弦查找表\\s\_sine\_fullscale\\(编译期固定)};
  \node[blk, right=of t1] (t2) {幅度 / 相位缩放\\$A$ 缩放 + CH2\\索引相位偏移};
  \node[blk, right=of t2] (t3) {DMA 缓冲区\\circular, 128 点\\DMA2\_CH3/CH4};
  \node[blk, right=of t3] (t4) {TIM6 TRGO 触发\\$f_s=f_{sig}\times128$};
  \node[io, right=of t4] (t5) {DAC 输出\\PA4 / PA5};
  \draw[fl] (t1)--(t2);
  \draw[fl] (t2)--(t3);
  \draw[fl] (t3)--(t4);
  \draw[fl] (t4)--(t5);
\end{tikzpicture}}
\caption{DAC 正弦波生成数据流图}
\label{fig:dac-flow}
\end{figure}

\phantomsection\label{sec:adc-sampling}
\subsectiontitle{5.2}{信号采集}

\textbf{ADC 输入信号来源：}在实际桌面用电监控场景下，ADC 的三路输入信号均来自 miniIO 副板的调理输出。市电电压经副板电压互感器隔离缩放后，由运放叠加 1.65 V 直流偏置输出 VL（0--3.3 V 范围内以 1.65 V 为中心的交流信号），送入 PC0/ADC1\_IN10。负载电流经副板电流互感器转换为小信号，由采样电阻和运放调理后输出 iL，送入 PC3/ADC1\_IN13。漏电电流由火线和零线共同穿过漏电互感器产生差动信号，经运放调理后输出 iLK，送入 PC2/ADC1\_IN12。市电频率由副板整形为低压方波 fr，送入 PA1/TIM2\_CH2 进行输入捕获。在调试阶段，也可将 DAC 输出通过跳线连接至 ADC 输入端进行环回测试，此时 ADC 采集的是 STM32 自身生成的正弦波信号。

ADC 采样硬件链路：TIM3（6400 Hz TRGO）$\rightarrow$ ADC1（扫描模式，3 通道）$\rightarrow$ DMA1\_Channel1（循环搬运）。

\textbf{采样配置：}ADC1 配置为独立模式、非连续扫描、TIM3 TRGO 外部触发。规则组 3 个 Rank 分别为 PC0/IN10（电压 VL）、PC3/IN13（电流 iL）、PC2/IN12（漏电电流 iLK），采样时间均取 239.5 ADC 时钟周期（约 20 $\mu$s @ 12 MHz），较长的采样窗口有利于降低外部信号处理电路输出阻抗的影响。

\textbf{DMA 传输：}DMA1\_Channel1 配置为外设到内存、半字宽度、循环模式，缓冲区大小为 $2 \times 3 \times 128 = 768$ 半字。HT（半传输）中断在前半帧 [0, 384) 写完时触发，TC（传输完成）中断在后半帧 [384, 768) 写完时触发。

\textbf{三层快照架构：}为解决 DMA ISR、主循环和 LCD/串口之间的数据竞争，采用三层缓冲设计：

第一层（ISR 写缓冲）：DMA HT/TC 中断中，deinterleave\_frame() 从 DMA 缓冲区将交错数据分离到双缓冲 s\_adc\_frame[2] 的写侧，完成后切换 write/ready 索引。

第二层（主循环工作帧）：app\_adc\_task() 在极短临界区（仅关中断取 ready 帧编号并清标志）后，将整帧数据从 ISR 缓冲复制到工作帧 s\_adc\_work\_frame（临界区外，不阻塞中断），再调用算法库计算电参量。

第三层（显示帧快照）：计算完成后将工作帧复制到 s\_adc\_display\_frame，LCD 波形绘制、串口 STATUS 输出和 CAL ZERO 校准统一读取此快照，保证数据一致性。

% [占位图]
\begin{figure}[H]
\centering
\begin{tikzpicture}[node distance=11mm, >=Latex]
  \node[blk, align=center] (isr) {第一层 · DMA HT/TC ISR（中断所有）\\deinterleave\_frame() 通道分离\\双缓冲 s\_adc\_frame[2]（写 / 就绪切换）};
  \node[blk, below=of isr, align=center] (work) {第二层 · 主循环工作帧（主循环所有）\\app\_adc\_task() 复制整帧 $\to$ s\_adc\_work\_frame\\$\to$ pm\_calc\_electrical 计算电参量};
  \node[blk, below=of work, align=center] (disp) {第三层 · 显示帧快照（共享只读）\\s\_adc\_display\_frame};
  \node[io, below=of disp, align=center] (cons) {LCD 波形绘制 · 串口 STATUS 输出 · CAL ZERO 校准};
  \draw[fl] (isr)-- node[lbl,right]{极短临界区：关中断取就绪帧号 + 清标志} (work);
  \draw[fl] (work)-- node[lbl,right]{复制整帧（临界区外，不阻塞中断）} (disp);
  \draw[fl] (disp)-- node[lbl,right]{只读快照} (cons);
\end{tikzpicture}
\caption{ADC 三层快照架构图}
\label{fig:adc-snapshot}
\end{figure}

\phantomsection\label{sec:adc-protocol}
\subsectiontitle{5.3}{通信协议}

与 ADC 采样和 DAC 输出相关的串口命令如表 \ref{tab:adc-commands} 所示。

\begin{table}[H]
\centering
\zihao{-5}
\caption{ADC/DAC 相关串口命令}
\label{tab:adc-commands}
\begin{tabular}{lll}
\toprule
命令 & 方向 & 说明 \\
\midrule
DAC SET MODE SINGLE\textbar DUAL & 主机$\rightarrow$设备 & 切换 DAC 工作模式 \\
DAC SET FREQ \textit{hz} & 主机$\rightarrow$设备 & 设置 DAC 输出频率（1--1000 Hz） \\
DAC SET AMP \textit{code} & 主机$\rightarrow$设备 & 设置 DAC 幅度码值（0--2047） \\
DAC SET PHASE \textit{deg} & 主机$\rightarrow$设备 & 设置 DAC CH2 相位偏移（0--359 度） \\
CAL ZERO & 主机$\rightarrow$设备 & 执行 ADC 零偏移校准 \\
OK DAC MODE=\ldots & 设备$\rightarrow$主机 & DAC 模式设置响应 \\
OK DAC FREQ=\ldots Hz & 设备$\rightarrow$主机 & DAC 频率设置响应 \\
OK DAC AMP=\ldots & 设备$\rightarrow$主机 & DAC 幅度设置响应 \\
OK DAC PHASE=\ldots & 设备$\rightarrow$主机 & DAC 相位设置响应 \\
OK CAL ZERO DONE & 设备$\rightarrow$主机 & 校准完成响应 \\
\bottomrule
\end{tabular}
\end{table}

STATUS? 响应中包含完整的 ADC/DAC 状态信息，典型格式为：

\begin{center}
\texttt{OK STATUS MEAS=50.00Hz PWM=50Hz REPORT=ON DAC\_MODE=DUAL DAC\_FREQ=50Hz DAC\_AMP=1500 DAC\_PHASE=0 VRMS=10.60 IRMS=1.060 ILK=0.394 P=11.2 S=11.2 PF=1.000 CAL=YES}
\end{center}

\textbf{DAC 参数修改流程：}收到 DAC SET 命令后，先读取当前配置（app\_dac\_get\_config），修改对应字段，调用 app\_dac\_apply\_config 执行：停止 DMA 和 TIM6 $\rightarrow$ 重建波形缓冲区 $\rightarrow$ 重新配置 DMA $\rightarrow$ 更新 TIM6 参数 $\rightarrow$ 重新启动。再读回配置（因钳位可能改变实际值）后格式化响应。

\textbf{CAL ZERO 校准：}取显示帧各通道 128 点算术平均值作为 DC 偏移补偿值。CAL ZERO 建议在无交流输入或输入稳定为静态中点时执行。若在带有交流信号、非整周期信号或偏置异常状态下执行，可能引入零点误差，导致后续 VRMS、IRMS、P 和 PF 计算产生偏差。

\phantomsection\label{sec:adc-waveform}
\subsectiontitle{5.4}{波形曲线绘制}

波形绘制在 LCD 的 DA Wave 页面完成，显示 DAC 双通道的 128 点波形预览。

\textbf{坐标映射：}12 位 DAC 样点（0--4095）映射到图形区域 y 坐标：

$$y = \text{top} + (4095 - \text{sample}) \times (\text{height} - 1) / 4095$$

其中图形区域 top = 132、height = 120 像素。样点值越大（电压越高），y 坐标越小（屏幕上方），符合直觉。

\textbf{绘制方法：}采用逐点连线方式。128 个样点均匀分布在图形宽度内，x 坐标为 $\text{left} + i \times (\text{width} - 1) / 127$。相邻两点之间调用 LCD\_DrawLine 绘制线段，形成连续波形曲线。CH1 使用绿色、CH2 使用黄色。

\textbf{坐标轴：}绘制纵轴（左侧灰线，含箭头）、横轴（零线，中线位置，含箭头）、1/4 和 3/4 处幅值刻度线、横轴等距时间刻度线。

\textbf{刷新策略：}LCD 刷新由两种条件触发——状态变化时即时刷新（display\_changed 标志）或固定 500 ms 周期刷新。换页时先清除内容区域避免旧页面残留。

% [占位图]
\begin{figure}[H]
\centering
\includegraphics[width=0.42\textwidth]{figures/lcd_dawave.png}
\caption{LCD DA Wave 页面显示效果（DUAL，50\,Hz，Amp 1500，相位 90°；CH1 绿、CH2 黄）}
\label{fig:lcd-da-wave}
\end{figure}

\clearpage

% ==============================
% 第 6 章：电参量计算与分析
% ==============================
\phantomsection\label{sec:calculation}
\sectiontitle{6}{电参量计算与分析}

电参量计算是系统的核心功能之一，将 ADC 采集的原始样点转换为有意义的电气物理量。本章介绍交流电参量计算方法、本系统的程序实现方案以及与上位机的通信设计。

\phantomsection\label{sec:calculation-method}
\subsectiontitle{6.1}{交流电参量计算方法}

交流电参量的计算方法主要有以下三类：

\textbf{（1）直接计算方法。}对采样得到的离散时域信号直接进行数学运算。电压有效值（RMS）定义为 $V_{RMS} = \sqrt{\frac{1}{N}\sum_{i=0}^{N-1}(V_i - V_{dc})^2}$，有功功率定义为 $P = \frac{1}{N}\sum_{i=0}^{N-1} V_i \times I_i$。该方法计算简单、实时性好，对采样率要求不高，适合嵌入式系统。

\textbf{（2）DFT计算方法。}将时域信号变换到频域，提取基波和各次谐波分量。可分别计算各谐波的幅值和相位，进而得到各次谐波功率。计算量为 $O(N^2)$，在样点数较多时开销较大。

\textbf{（3）FFT计算方法。}DFT 的快速算法，计算量 $O(N \log N)$。适合样点数较多（通常要求 2 的幂次）且需要频谱分析的场景。在 Cortex-M3 上实现 FFT 需要额外的整型数蝶形运算库。

系统选择直接计算方法：每周期 128 点采样足以满足基波 RMS 与功率精度，计算量小、实时性好，且无需复数运算与频谱分析；Cortex-M3 无硬件 FPU，整数运算效率更高。

\phantomsection\label{sec:calculation-program}
\subsectiontitle{6.2}{程序实现}

电参量计算由 pm\_calc\_electrical() 函数完成，以 no\_std 整型方式实现，通过 C 接口供 C 代码调用。整数运算无运行时浮点开销，契合 Cortex-M3 无 FPU 的硬件特点。计算流程如下：

\textbf{（1）DC 偏移确定。}各通道独立计算直流分量。若已通过 CAL ZERO 命令手动校准（provided\_zero $\neq$ 0），直接使用校准值；否则从本批 128 个样点取算术平均值 $Z = \frac{1}{N}\sum_{i=0}^{N-1} x_i（N=128）$ 作为自动估算的 DC 偏移。

\textbf{（2）RMS 码值计算。}对每通道去直流后计算均方根值：

$$\text{RMS}_{code} = \sqrt{\frac{1}{N}\sum_{i=0}^{N-1}(x_i - Z)^2}$$

中间平方和使用 u64（64 位无符号整数）累加，12 位 ADC 码值的平方最大 $2047^2 \approx 4.19 \times 10^6$，128 点累加最大约 $5.37 \times 10^8$，远小于 u64 上限 $1.84 \times 10^{19}$，不会溢出。开平方运算采用高效的整数开平方算法实现，返回所求平方根的整数部分。

\textbf{（3）物理量换算。}RMS 码值通过校准增益转换为物理量：

$$V_{rms} \times 100 = \text{RMS}_{code} \times G_v / 1000$$
$$I_{rms} \times 1000 = \text{RMS}_{code} \times G_i / 1000$$

其中 $G_v$、$G_i$ 为校准增益$\times$1000（默认 1000，即 1:1）。

\textbf{（4）有功功率计算。}逐点计算去直流后的瞬时功率积并取均值：

$$P_{code} = \frac{1}{N}\sum_{i=0}^{N-1}(v_i - Z_v) \times (i_i - Z_i)$$

有功功率物理量：

$$P \times 10 = P_{code} \times G_v \times G_i / 10^{10}$$

除数为 $10^{10}$ 的原因是 $V_{rms}\times100$（单位 0.01 V）和 $I_{rms}\times1000$（单位 0.001 A）的单位缩放因子不同，$\times$增益$^2$后需除以增益$^2 \times 10000$。结果保留符号，正值表示消耗功率，负值表示回馈功率。

\textbf{（5）视在功率：}$S \times 10 = V_{rms\times100} \times I_{rms\times1000} / 10000$。

\textbf{（6）功率因数：}$\text{PF} \times 1000 = |P| \times 1000 / S$，钳位到 [0, 1000] 区间。$S = 0$ 时（无负载）PF 返回 0，因为无负载时功率因数无物理意义。

% [占位图]
\begin{figure}[H]
\centering
\begin{tikzpicture}[node distance=7mm and 12mm, >=Latex]
  \node[blk] (p1) {ADC 128 点样点（三通道）};
  \node[dec, below=of p1] (d) {已 CAL ZERO?};
  \node[blk, left=of d, align=center] (a0) {取本批均值\\$Z=\sum_{i=0}^{N-1} x_i/128$};
  \node[blk, right=of d, align=center] (a1) {使用校准\\DC 偏移};
  \node[blk, below=16mm of d] (p2) {去直流 RMS 码值\\$\text{RMS}=\sqrt{\tfrac1N\sum_{i=0}^{N-1}(x_i-Z)^2}$（整数开方）};
  \node[blk, below=of p2] (p3) {物理量换算\\$V_{rms}{\times}100=\text{RMS}\cdot G_v/1000$\\$I_{rms}{\times}1000=\text{RMS}\cdot G_i/1000$};
  \node[blk, below=of p3] (p4) {有功功率\\$P{\times}10=\overline{(v{-}Z_v)(i{-}Z_i)}\cdot G_vG_i/10^{10}$};
  \node[blk, below=of p4] (p5) {视在功率\\$S{\times}10=V_{rms\times100}\cdot I_{rms\times1000}/10^{4}$};
  \node[blk, below=of p5] (p6) {功率因数\\$\text{PF}{\times}1000=|P|\cdot1000/S$，钳位 $[0,1000]$};
  \draw[fl] (p1)--(d);
  \draw[fl] (d)-- node[lbl,above]{否} (a0);
  \draw[fl] (d)-- node[lbl,above]{是} (a1);
  \draw[fl] (a0) |- (p2.west);
  \draw[fl] (a1) |- (p2.east);
  \draw[fl] (p2)--(p3);
  \draw[fl] (p3)--(p4);
  \draw[fl] (p4)--(p5);
  \draw[fl] (p5)--(p6);
\end{tikzpicture}
\caption{电参量计算流程图}
\label{fig:calc-flow}
\end{figure}

\phantomsection\label{sec:calculation-comm}
\subsectiontitle{6.3}{通信设计}

电参量通过串口协议的 STATUS? 命令和自动上报功能传输至上位机。

\textbf{STATUS? 查询：}主机发送 STATUS? 命令帧，设备读取当前监控快照（app\_monitor\_state\_read），调用 app\_command\_format\_status 将所有电参量格式化为一行文本，封装为 TYPE=0x81 响应帧返回。格式化使用 snprintf，各参数的小数位数格式化统一：VRMS 两位小数（0.01 V）、IRMS 和 ILK 三位小数（0.001 A）、P 一位小数（0.1 W，带符号）、S 一位小数（0.1 VA）、PF 三位小数、CAL 显示 YES/NO。

\textbf{自动上报：}开启 REPORT ON 后，主循环每秒检查一次 take\_report 信号（由 app\_capture\_task 按上报周期设置），触发时读取监控快照并格式化 STATUS 行，封装为 TYPE=0x82 事件帧发送。上位机侧无需逐条发送查询命令即可持续接收设备状态更新。

\textbf{上位机工具：}项目配套的串口工具（tools/serial\_cmd/），支持命令行交互和协议帧解析。上位机发送文本命令，工具自动封装为二进制帧并计算 CRC；接收到响应帧后解封装并显示 Payload 文本。

% [占位图]
\begin{figure}[H]
\centering
\begin{lstlisting}[style=serialterm]
$ serial_cmd /dev/cu.usbserial-110 "STATUS?"
>> STATUS?
<< RSP OK STATUS MEAS=50.00Hz PWM=50Hz REPORT=ON DAC_MODE=DUAL DAC_FREQ=50Hz DAC_AMP=1500 DAC_PHASE=0 VRMS=10.60 IRMS=1.060 ILK=0.394 P=11.2 S=11.2 PF=1.000 CAL=YES

$ serial_cmd /dev/cu.usbserial-110 "REPORT ON"
>> REPORT ON
<< RSP OK REPORT=ON
--- 继续监听 5 秒自动上报 ---
<< EVT OK STATUS MEAS=50.00Hz PWM=50Hz REPORT=ON DAC_MODE=DUAL DAC_FREQ=50Hz DAC_AMP=1500 DAC_PHASE=0 VRMS=10.60 IRMS=1.060 ILK=0.394 P=11.2 S=11.2 PF=1.000 CAL=YES
<< EVT OK STATUS MEAS=50.00Hz PWM=50Hz REPORT=ON DAC_MODE=DUAL DAC_FREQ=50Hz DAC_AMP=1500 DAC_PHASE=0 VRMS=10.61 IRMS=1.061 ILK=0.394 P=11.2 S=11.2 PF=1.000 CAL=YES
\end{lstlisting}
\caption{串口工具 STATUS? 查询与自动上报}
\label{fig:serial-status}
\end{figure}

\clearpage

% ==============================
% 第 7 章：系统调试
% ==============================
\phantomsection\label{sec:debug}
\sectiontitle{7}{系统调试}

本章对系统的串口通信协议、PWM 频率输出、输入捕获测频、ADC 采样、电参量计算以及 miniIO 联机运行进行调试。调试过程采用"命令下发——设备响应——LCD 显示——数据记录——结果分析"的方式进行，重点验证通信协议可靠性、频率测量正确性、ADC 采样稳定性和电参量计算趋势正确性。

\phantomsection\label{sec:debug-protocol}
\subsectiontitle{7.1}{通信协议分析}

系统通信协议采用两层设计，表 \ref{tab:protocol-summary} 汇总了各帧类型及其用途。

\begin{table}[H]
\centering
\zihao{-5}
\caption{协议帧类型汇总}
\label{tab:protocol-summary}
\begin{tabular}{llll}
\toprule
帧类型 & TYPE 值 & 方向 & 用途 \\
\midrule
命令帧 & 0x01 & 主机$\rightarrow$设备 & 发送控制命令 \\
响应帧 & 0x81 & 设备$\rightarrow$主机 & 命令执行结果响应 \\
事件帧 & 0x82 & 设备$\rightarrow$主机 & 自动上报 / 启动就绪通知 \\
错误帧 & 0x83 & 设备$\rightarrow$主机 & 帧格式错误通知 \\
\bottomrule
\end{tabular}
\end{table}

\textbf{帧格式详解：}

\begin{table}[H]
\centering
\zihao{-5}
\caption{二进制帧结构}
\label{tab:frame-struct}
\begin{tabular}{lcccl}
\toprule
字段 & 偏移 & 长度 & 字节序 & 说明 \\
\midrule
SOF0 & 0 & 1 & -- & 固定 0xA5 \\
SOF1 & 1 & 1 & -- & 固定 0x5A \\
VERSION & 2 & 1 & -- & 协议版本 0x01 \\
TYPE & 3 & 1 & -- & 帧类型 \\
SEQ & 4 & 1 & -- & 序号 \\
LEN & 5--6 & 2 & 小端 & Payload 长度 \\
PAYLOAD & 7--(6+LEN) & LEN & -- & 文本载荷 \\
CRC16 & 末尾 2 字节 & 2 & 小端 & CRC-16/Modbus \\
\bottomrule
\end{tabular}
\end{table}

\textbf{CRC 校验范围：}从 VERSION 到 PAYLOAD 末尾（不含 SOF 和 CRC 本身），初始值 0xFFFF，多项式 0xA001（反转表示）。

\textbf{错误处理：}CRC 校验失败时静默丢弃帧（不响应），版本号或帧类型不匹配时发送 TYPE=0x83 错误帧（payload ``ERR BAD\_FRAME''），Payload 超长（$>$ 63 字节）时同样返回错误帧。环形缓冲区溢出时重置解析器状态。

通信协议调试主要验证二进制帧的帧头识别、长度解析、Payload 承载、CRC 校验和错误处理功能。上位机通过串口工具发送 HELP、STATUS?、REPORT ON、PWM SET、DAC SET 等命令，设备端解析后返回响应帧。协议测试内容如表 \ref{tab:protocol-test} 所示。

\begin{table}[H]
\centering
\zihao{-5}
\caption{通信协议测试记录表}
\label{tab:protocol-test}
\begin{tabular}{lllll}
\toprule
测试项目 & 发送命令 & 期望结果 & 实际结果 & 结论 \\
\midrule
帮助命令 & HELP & 返回支持的命令列表 & 返回 4 行命令列表 & 通过 \\
状态查询 & STATUS? & 返回 VRMS、IRMS、P 等状态量 & 单帧返回 VRMS/IRMS/P 等 & 通过 \\
自动上报 & REPORT ON & 周期返回 TYPE=0x82 事件帧 & 每约 1 s 收到一帧 0x82 & 通过 \\
PWM 设置 & PWM SET 50 & PWM 频率更新为 50 Hz & 应答 OK PWM FREQ=50Hz & 通过 \\
DAC 设置 & DAC SET FREQ 50 & DAC 频率更新为 50 Hz & 应答 OK DAC FREQ=50Hz & 通过 \\
CRC 错误 & 修改 CRC 字节 & 设备丢弃或返回错误帧 & 翻转末字节后无响应 & 通过 \\
\bottomrule
\end{tabular}
\end{table}

通过测试可见，协议解析器能够正确识别合法帧，并能对异常帧进行过滤或错误响应，满足系统远程参数设置和状态上报需求。

% [占位图]
\begin{figure}[H]
\centering
\begin{lstlisting}[style=serialterm]
$ serial_cmd /dev/cu.usbserial-110 "STATUS?"
>> STATUS?
<< RSP OK STATUS MEAS=50.00Hz PWM=50Hz REPORT=ON DAC_MODE=DUAL DAC_FREQ=50Hz DAC_AMP=1500 DAC_PHASE=0 VRMS=10.60 IRMS=1.060 ILK=0.394 P=11.2 S=11.2 PF=1.000 CAL=YES

# STATUS? 命令帧（A5 5A 帧头 + CRC-16/Modbus，SEQ=0x2A 示例）
#   A5 5A(SOF) 01(VER) 01(TYPE) 2A(SEQ) 07 00(LEN) 53 54 41 54 55 53 3F("STATUS?") B4 59(CRC)
#   CRC 覆盖 VER..Payload，即 01 01 2A 07 00 53..3F，计算得 B4 59

$ crc_test /dev/cu.usbserial-110
=== 测试 1：正确 CRC 帧 (STATUS?)              收到 1 条响应 -> 正常响应
=== 测试 2：破坏 CRC 帧 (翻转 CRC 末字节)        收到 0 条响应 -> 被正确拒绝（静默丢弃）
=== 测试 3：篡改 payload 保留原 CRC (HELP->HALP)  收到 0 条响应 -> 被正确拒绝
\end{lstlisting}
\caption{串口协议帧与 CRC 校验}
\label{fig:protocol-hex}
\end{figure}

\phantomsection\label{sec:debug-frequency}
\subsectiontitle{7.2}{频率调试}

频率调试采用 PA8/TIM1\_CH1 输出 PWM 方波，PA1/TIM2\_CH2 输入捕获测量频率。调试时将 PA8 与 PA1 短接进行内部环回自测，或使用外部信号源输入 PA1。串口发送 PWM SET 命令修改输出频率，观察串口 STATUS? 返回的 PWM 设定值和 MEAS 测量值，同时记录 LCD 的 Freq Measure 页面显示结果。

\begin{table}[H]
\centering
\zihao{-5}
\caption{频率调试记录表}
\label{tab:pwm-test}
\begin{tabular}{llllllll}
\toprule
目标频率 & 设置命令 & PSC & ARR & 仪器实测 & TIM2 测量值 & 相对误差/\% & 结论 \\
\midrule
50 Hz & PWM SET 50 & 21 & 65454 & 50.0011 Hz & 50.00 Hz & 0.00 & 通过 \\
100 Hz & PWM SET 100 & 10 & 65454 & 100.002 Hz & 100.00 Hz & 0.00 & 通过 \\
1 kHz & PWM SET 1000 & 1 & 35999 & 1000.03 Hz & 1000.00 Hz & 0.00 & 通过 \\
10 kHz & PWM SET 10000 & 0 & 7199 & 10000.3 Hz & 10000.00 Hz & 0.00 & 通过 \\
100 kHz & PWM SET 100000 & 0 & 719 & 100003 Hz & 100000.00 Hz & 0.00 & 通过 \\
\bottomrule
\end{tabular}
\end{table}

表 \ref{tab:pwm-test} 中"仪器实测"栏为 SIGLENT 数字示波器频率计的实测读数，作为外部基准与 TIM2 输入捕获测量值对比。由于 PA8 输出与 PA1 捕获共用 STM32 的同一 72 MHz 时钟，TIM2 测量值与目标频率完全吻合；示波器采用独立时基，其读数与目标的偏差均在 $\pm$0.003\% 以内，二者一致表明频率测量功能正确。

\textbf{超时验证：}断开 PA1 输入信号，观察 2 秒后 LCD 频率显示是否归零。

由测试结果可知，在测试频率范围内，TIM2 输入捕获能够正确反映外部方波频率。TIM2 采用 1 μs 计数分辨率，低频时相对分辨率高。本测试中 PWM 输出与输入捕获共用同一 72 MHz 时钟且分频均为整数，故实测各频点 TIM2 测量值与目标完全一致。对外部非同源信号，高频端单周期测量对计数量化更敏感；如需更高精度，可采用多周期平均法。PWM 频率的软件设定范围为 1--100000 Hz，实际输出波形质量受定时器 PSC/ARR 整数分频、IO 负载、PA8 板载 LED 连接以及测试设备带宽影响，高频端的真实分辨率限制需结合独立时基的示波器或逻辑分析仪实测结果判断。

\begin{figure}[H]
\centering
\begin{subfigure}[t]{0.48\textwidth}
  \centering
  \includegraphics[width=\linewidth]{figures/freq_50hz.png}
  \caption{50 Hz（实测 50.0011 Hz）}
\end{subfigure}
\hfill
\begin{subfigure}[t]{0.48\textwidth}
  \centering
  \includegraphics[width=\linewidth]{figures/freq_100hz.png}
  \caption{100 Hz（实测 100.002 Hz）}
\end{subfigure}

\vspace{0.3cm}
\begin{subfigure}[t]{0.48\textwidth}
  \centering
  \includegraphics[width=\linewidth]{figures/freq_1khz.png}
  \caption{1 kHz（实测 1.00003 kHz）}
\end{subfigure}
\hfill
\begin{subfigure}[t]{0.48\textwidth}
  \centering
  \includegraphics[width=\linewidth]{figures/freq_10khz.png}
  \caption{10 kHz（实测 10.0003 kHz）}
\end{subfigure}

\vspace{0.3cm}
\begin{subfigure}[t]{0.48\textwidth}
  \centering
  \includegraphics[width=\linewidth]{figures/freq_100khz.png}
  \caption{100 kHz（实测 100.003 kHz）}
\end{subfigure}

\caption{各频率点示波器实测波形（PA8 输出，PA8--PA1 短接环回）}
\label{fig:debug-freq}
\end{figure}

\phantomsection\label{sec:debug-ad}
\subsectiontitle{7.3}{AD采样调试}

ADC 采样调试分为零输入基线测试、DAC 环回测试和波形显示测试。零输入测试时不接外部交流信号，观察 PC0、PC3、PC2 三个 ADC 通道的原始码值是否稳定在中点附近。DAC 环回测试时，将 PA4/DAC\_OUT1 和 PA5/DAC\_OUT2 通过跳线接入 ADC 输入端，设置不同幅值、频率和相位，观察采样结果和电参量变化。

\textbf{（1）零信号基线测试。}不接入外部信号，发送 STATUS? 查看各通道 ADC 原始码值范围。PC0 和 PC3 在无信号时应稳定在中点附近（约 2048），若某通道码值明显偏大或偏小，说明引脚被外部电路拉偏。

\textbf{（2）DAC 环回测试。}将 DAC 输出（PA4/PA5）通过跳线连接到 ADC 输入（PC0/PC3），发送 DAC SET MODE DUAL、DAC SET FREQ 50、DAC SET AMP 1500 等命令配置测试信号，观察 ADC 采样波形和电参量。环回测试前应确认 CAL ZERO 的执行条件。CAL ZERO 用于记录三通道 ADC 的直流零点，建议在无交流输入或输入稳定为静态中点时执行。若在带有交流信号、非整周期信号或偏置异常状态下执行，可能引入零点误差，导致后续 VRMS、IRMS、P 和 PF 计算产生偏差。

\textbf{（3）波形验证。}在 LCD DA Wave 页面观察 DAC 波形预览是否为标准正弦波；在串口查看 VRMS 和 IRMS 值是否与 DAC 幅度设置一致。

\begin{table}[H]
\centering
\zihao{-5}
\caption{ADC 原始采样记录表}
\label{tab:adc-raw}
\begin{tabular}{lllllll}
\toprule
测试条件 & 通道 & 最小值 & 最大值 & 平均值 & 峰峰值 & 结论 \\
\midrule
零输入 & PC0/VL & 待填 & 待填 & 待填 & 待填 & 待填 \\
零输入 & PC3/iL & 待填 & 待填 & 待填 & 待填 & 待填 \\
零输入 & PC2/iLK & 待填 & 待填 & 待填 & 待填 & 待填 \\
DAC 50 Hz AMP=1500 & PC0/VL & 待填 & 待填 & 待填 & 待填 & 待填 \\
DAC 50 Hz AMP=1500 & PC3/iL & 待填 & 待填 & 待填 & 待填 & 待填 \\
\bottomrule
\end{tabular}
\end{table}

ADC 在 TIM3 触发下按 6400 Hz 完成三通道扫描，DMA 半传输与传输完成中断分别分离前后半帧的通道数据；该链路的实际效果在 7.4 节电参量调试中结合 VRMS、IRMS 随设定信号的变化加以验证。

% [占位图]
\begin{figure}[H]
\centering
\begin{lstlisting}[style=serialterm]
# DAC->ADC 跳线环回 (PA4->PC0, PA5->PC3)
$ serial_cmd /dev/cu.usbserial-110 "DAC SET MODE DUAL"
>> DAC SET MODE DUAL
<< RSP OK DAC MODE=DUAL
$ serial_cmd /dev/cu.usbserial-110 "DAC SET FREQ 50"
>> DAC SET FREQ 50
<< RSP OK DAC FREQ=50Hz
$ serial_cmd /dev/cu.usbserial-110 "DAC SET AMP 1500"
>> DAC SET AMP 1500
<< RSP OK DAC AMP=1500
$ serial_cmd /dev/cu.usbserial-110 "STATUS?"
>> STATUS?
<< RSP OK STATUS MEAS=50.00Hz PWM=50Hz REPORT=OFF DAC_MODE=DUAL DAC_FREQ=50Hz DAC_AMP=1500 DAC_PHASE=0 VRMS=10.60 IRMS=1.060 ILK=0.394 P=11.2 S=11.2 PF=1.000 CAL=YES
# LCD AC Measure 页同步显示相同读数及电压通道原始码范围 (V ADC min/max)
\end{lstlisting}
\caption{ADC 采样调试结果}
\label{fig:debug-adc}
\end{figure}

\phantomsection\label{sec:debug-calculation}
\subsectiontitle{7.4}{电参量调试}

电参量调试主要验证 VRMS、IRMS、有功功率 P、视在功率 S 和功率因数 PF 的计算趋势。由于本系统未使用标准功率源进行标定，本节重点验证相对变化关系和计算逻辑，不将测试结果表述为计量级精度。

\textbf{（1）幅值线性验证。}DAC 频率固定 50 Hz、相位差 0°，逐步改变 DAC 幅值（500、1000、1500、2000），记录 VRMS、IRMS、P、S、PF，验证 VRMS 是否随幅值近似线性增长。

\begin{table}[H]
\centering
\zihao{-5}
\caption{幅值线性测试记录表}
\label{tab:amp-linear}
\begin{tabular}{lllllllll}
\toprule
DAC 频率 & 幅值码 & 相位差 & VRMS/V & IRMS/A & P/W & S/VA & PF & 现象分析 \\
\midrule
50 Hz & 500 & 0° & 3.53 & 0.353 & 1.2 & 1.2 & 1.000 & 同相 PF=1，P=S \\
50 Hz & 1000 & 0° & 7.07 & 0.706 & 4.9 & 4.9 & 1.000 & V、I 约 2×，P 约 4× \\
50 Hz & 1500 & 0° & 10.60 & 1.060 & 11.2 & 11.2 & 1.000 & V、I 约 3×，P 约 9× \\
50 Hz & 2000 & 0° & 14.11 & 1.399 & 19.7 & 19.7 & 1.000 & V、I 约 4×，P 约 16× \\
\bottomrule
\end{tabular}
\end{table}

\textbf{（2）相位验证。}DAC 频率 50 Hz、幅值 1500，分别设置相位差为 0°、90°、180°，观察 P 的符号和 PF 的变化趋势。同相时 P $>$ 0（消耗功率），反相时 P $<$ 0（回馈功率），正交时 PF 应接近 0。

\begin{table}[H]
\centering
\zihao{-5}
\caption{相位测试记录表}
\label{tab:phase-test}
\begin{tabular}{llllll}
\toprule
DAC 频率 & 幅值码 & 相位差 & P 符号 & PF 理论趋势 & PF 实测值 \\
\midrule
50 Hz & 1500 & 0° & 正 & 接近 1 & 1.000 \\
50 Hz & 1500 & 90° & 接近 0 & 接近 0 & 0.044 \\
50 Hz & 1500 & 180° & 负 & 接近 1 & 1.000 \\
\bottomrule
\end{tabular}
\end{table}

\textbf{（3）无负载情况。}无信号输入时 S = 0，PF 应返回 0（非 NaN 或无穷大）。

\textbf{（4）CAL ZERO 校准验证。}无信号时发送 CAL ZERO，再发送 STATUS? 确认 CAL=YES，各 RMS 值应接近 0。

测试结果显示，幅值增大时 VRMS 和 IRMS 随之增大；同相输入时有功功率为正，反相输入时有功功率为负；正交输入时功率因数接近 0，符合交流电参量计算规律。

% [占位图]
\begin{figure}[H]
\centering
\begin{lstlisting}[style=serialterm]
# 幅值线性（DUAL，50 Hz，相位 0），STATUS? 省略前缀字段
>> DAC SET AMP 1000
<< RSP OK DAC AMP=1000
>> STATUS?
<< RSP OK STATUS ... DAC_AMP=1000 DAC_PHASE=0 VRMS=7.07 IRMS=0.706 P=4.9 S=4.9 PF=1.000 CAL=YES
>> DAC SET AMP 1500
<< RSP OK DAC AMP=1500
>> STATUS?
<< RSP OK STATUS ... DAC_AMP=1500 DAC_PHASE=0 VRMS=10.60 IRMS=1.060 P=11.2 S=11.2 PF=1.000 CAL=YES

# 相位（DUAL，50 Hz，幅值 1500）
>> DAC SET PHASE 90
<< RSP OK DAC PHASE=90
>> STATUS?
<< RSP OK STATUS ... DAC_AMP=1500 DAC_PHASE=90 VRMS=10.60 IRMS=1.060 P=0.5 S=11.2 PF=0.044 CAL=YES
>> DAC SET PHASE 180
<< RSP OK DAC PHASE=180
>> STATUS?
<< RSP OK STATUS ... DAC_AMP=1500 DAC_PHASE=180 VRMS=10.60 IRMS=1.060 P=-11.2 S=11.2 PF=1.000 CAL=YES
\end{lstlisting}
\caption{电参量调试结果}
\label{fig:debug-calc}
\end{figure}

\phantomsection\label{sec:debug-online}
\subsectiontitle{7.5}{联机调试}

联机调试将 MiniSTM32 开发板通过串口连接上位机，同时观察 LCD 显示和串口数据。调试分为 DAC 环回测试和 miniIO 实际接入测试两部分。MiniSTM32 开发板通过 CH340 USB 转串口芯片（PA9/TX $\rightarrow$ CH340 RX，PA10/RX $\rightarrow$ CH340 TX）连接上位机 USB 口，串口配置 115200 8N1。

\textbf{（1）DAC 环回测试。}将 PA4/DAC\_OUT1、PA5/DAC\_OUT2 通过跳线连接至 ADC 输入端，用于验证 ADC 采样、DMA 数据拆分、电参量计算和串口上报功能。该模式不接入市电，便于调整正弦波幅值、频率和相位，并观察 VRMS、IRMS、P、S、PF 等参数是否随输入信号变化。

\textbf{（2）miniIO 实际接入测试。}MiniSTM32 与 miniIO 通过 JP1 排线连接，5V 和 GND 为副板供电，VL 接 PC0/ADC1\_IN10，iL 接 PC3/ADC1\_IN13，iLK 接 PC2/ADC1\_IN12，fr 接 PA1/TIM2\_CH2。负载接入 miniIO 的交流输入输出端后，系统通过 ADC 采集电压、电流和漏电电流信号，通过 TIM2 输入捕获测量市电频率，并在 LCD 与串口 STATUS 数据中显示计算结果。联机调试时需同时记录 miniSTM32 与 miniIO 接线图、LCD 显示页面、串口助手返回数据和若干组稳定性测试数据。

\begin{table}[H]
\centering
\zihao{-5}
\caption{miniSTM32 与 miniIO 接线记录表}
\label{tab:miniio-wiring}
\begin{tabular}{llll}
\toprule
miniIO 信号 & STM32 连接 & 功能 & 检查结果 \\
\midrule
5V & MiniSTM32 5V & 副板供电 & 待填 \\
GND & MiniSTM32 GND & 共地 & 待填 \\
VL & PC0/ADC1\_IN10 & 电压检测 & 待填 \\
iL & PC3/ADC1\_IN13 & 电流检测 & 待填 \\
iLK & PC2/ADC1\_IN12 & 漏电电流检测 & 待填 \\
fr & PA1/TIM2\_CH2 & 频率测量 & 待填 \\
\bottomrule
\end{tabular}
\end{table}

\textbf{联机验证内容：}

（1）上电后串口收到事件帧 ``OK COURSE1 READY''，确认固件启动正常。

（2）LCD 显示 Main Menu 页面，通过 KEYUP/KEY1/KEY0 导航各子页面，确认 LCD 渲染正确。

（3）通过串口发送各命令，验证响应帧内容正确、CRC 校验通过。

（4）开启 REPORT ON，观察自动上报数据流是否按约 1 秒周期稳定推送。

（5）对比 LCD 显示数值和串口 STATUS? 查询结果，确认两端数据一致。

联机运行后，串口发送 STATUS? 获取系统状态，并与 LCD 页面显示结果对比。若两端 VRMS、IRMS、P、S、PF 和频率显示一致，说明 LCD 和串口均读取同一监控快照，数据链路一致。开启 REPORT ON 后，连续观察自动上报数据，记录是否存在丢帧、乱码、卡死或显示异常。

\begin{table}[H]
\centering
\zihao{-5}
\caption{联机稳定性测试记录表}
\label{tab:stability-test}
\begin{tabular}{llllll}
\toprule
测试时间 & 上报周期 & 串口异常次数 & LCD 刷新 & 数据波动范围 & 结论 \\
\midrule
1 min & 1 s & 0 & 正常 & VRMS 10.60--10.61 V & 稳定，无丢帧 \\
3 min & 1 s & 0 & 正常 & VRMS 10.60--10.61 V & 稳定，无丢帧 \\
5 min & 1 s & 0 & 正常 & VRMS 10.60--10.61 V & 稳定，无丢帧 \\
\bottomrule
\end{tabular}
\end{table}

% [占位图]
\begin{figure}[H]
\centering
\fbox{\parbox{0.8\textwidth}{\centering\vspace{3cm}
\textit{图占位：miniSTM32 与 miniIO 联机接线图}\\[0.3cm]
建议绘制标注清楚的连接示意图：MiniSTM32 开发板 JP1 排针各信号（5V/GND/VL/iL/iLK/fr/beep/relay）与 miniIO 副板对应接口的连线关系，标注引脚编号和信号名称。不要只放全景照片，必须有标注清楚的接线图。
\vspace{3cm}}}
\caption{miniSTM32 与 miniIO 联机接线图}
\label{fig:debug-miniio-wiring}
\end{figure}

% [占位图]
\begin{figure}[H]
\centering
\fbox{\parbox{0.8\textwidth}{\centering\vspace{3cm}
\textit{图占位：联机调试全景照片}\\[0.3cm]
建议拍摄包含 MiniSTM32 开发板、miniIO 副板、LCD 屏幕显示、USB 线缆连接、负载接入、上位机串口助手窗口的全景照片。
\vspace{3cm}}}
\caption{联机调试全景}
\label{fig:debug-online}
\end{figure}

% [占位图]
\begin{figure}[H]
\centering
\begin{subfigure}[t]{0.30\textwidth}
  \centering
  \includegraphics[width=\linewidth]{figures/lcd_menu.png}
  \caption{Main Menu}
\end{subfigure}
\hfill
\begin{subfigure}[t]{0.30\textwidth}
  \centering
  \includegraphics[width=\linewidth]{figures/lcd_square.png}
  \caption{Square Wave}
\end{subfigure}
\hfill
\begin{subfigure}[t]{0.30\textwidth}
  \centering
  \includegraphics[width=\linewidth]{figures/lcd_measure.png}
  \caption{Freq Measure}
\end{subfigure}

\vspace{0.3cm}
\begin{subfigure}[t]{0.30\textwidth}
  \centering
  \includegraphics[width=\linewidth]{figures/lcd_dawave.png}
  \caption{DA Wave}
\end{subfigure}
\hspace{0.03\textwidth}
\begin{subfigure}[t]{0.30\textwidth}
  \centering
  \includegraphics[width=\linewidth]{figures/lcd_acmeasure.png}
  \caption{AC Measure}
\end{subfigure}
\caption{LCD 各页面显示效果}
\label{fig:lcd-pages}
\end{figure}

\clearpage

% ==============================
% 第 8 章：总结与展望
% ==============================
\phantomsection\label{sec:summary}
\sectiontitle{8}{总结与展望}

本次课程设计完成了基于 STM32F103RCT6 的桌面用电监控系统：信号生成、频率测量、DAC 输出、ADC 采样、电参量计算、LCD 显示与串口通信各模块均已实现并联调通过。系统采用裸机协作调度，无 RTOS 依赖，软件分层清晰，LCD 与串口通过统一监控快照保持数据一致。

\textbf{调试过程中遇到的主要问题及解决办法：}

（1）PC1（ADC1\_IN11）电流通道读数异常。现象：DAC DUAL 环回测试时 IRMS 始终约为 0.056，远小于预期。排查发现 MiniSTM32 开发板的 PC1 与板载触摸屏 T\_PEN 共享引脚，硬件上拉电阻将 PC1 拉至约 3V，导致 ADC 读数严重偏高。解决办法：将电流采样通道迁移至 PC3（ADC1\_IN13），避开触摸屏干扰。

（2）ADC 采样数据竞争。现象：LCD 显示的电参量偶尔跳变，串口 STATUS 返回的 ADC 样点范围不一致。根因是 DMA ISR 正在写入缓冲区时主循环同时在读取。解决办法：设计三层快照架构——ISR 写双缓冲、主循环极短临界区取帧号、计算后复制到显示帧——从根本上消除数据竞争。

（3）命令响应栈溢出。现象：增大 RESPONSE\_MAX（响应行数）后设备硬故障无响应。根因是 app\_command\_result\_t 结构体在 dispatch\_command\_frame() 栈上分配，4$\times$256 = 1024 字节已接近栈上限。解决办法：将 RESPONSE\_MAX 限制为 4，并将 HELP 命令的两条 DAC 帮助合并为一行以适配行数限制。

（4）CAL ZERO 执行时机不当导致电参量异常。现象：在 DAC 输出交流信号时执行 CAL ZERO 后，VRMS/IRMS/P 等参数出现明显偏差。根因是校准函数在带有交流信号的状态下记录了非静态零点，引入零点误差，导致后续 RMS 和功率计算偏差。解决办法：在文档和帮助信息中明确 CAL ZERO 应在无交流输入或输入稳定为静态中点时执行。

\textbf{心得体会：}

本次课程设计中最具挑战性的部分是数据竞争问题的排查——DMA 中断与主循环并发访问同一缓冲区的问题在代码审查中难以发现，需要结合现象反复推理数据时序。三层快照架构的设计过程加深了我对嵌入式并发数据一致性的理解。

在团队协作方面，与同学在引脚分配、协议格式和算法选择上进行了充分讨论，把常量集中到 app\_config.h、把帧定义放进共享的 app\_protocol\_defs.h，使固件与上位机工具的协议保持一致。AI 工具在代码审查、Bug 分析和架构建议方面提供了有价值的辅助，尤其是在数据竞争根因分析和栈溢出问题定位中发挥了关键作用；但 AI 生成的代码仍需人工审查边界条件和嵌入式约束。

\textbf{改进方向：}

（1）增加 FFT 频谱分析功能，支持谐波检测和 THD 计算。

（2）将串口协议升级为支持批量数据传输，便于绘制实时波形。

（3）增加 SD 卡数据记录功能，支持长时间监测数据离线分析。

（4）优化 LCD 显示刷新策略，采用局部刷新减少闪烁。

\end{document}
